\section{Obiettivi della sperimentazione}


La fase sperimentale ha l'obiettivo di valutare il comportamento del modello di ottimizzazione bi-obiettivo e delle due euristiche implementate, Greedy e Clarke-Wright, in differenti condizioni operative e configurazioni del problema. L'analisi è finalizzata a verificare come la soluzione ottenuta vari al modificarsi delle caratteristiche dell'istanza e dei parametri che definiscono il compromesso tra la qualità del servizio offerto ai cittadini e l'efficienza logistica del sistema di raccolta.

In particolare, la sperimentazione considera differenti dimensioni e configurazioni del problema, analizzando la variazione del numero di utenti attivi, la tipologia urbanistica e la distribuzione geografica degli utenti. A queste istanze sintetiche si affianca infine un'analisi condotta su una configurazione basata sulla mappa reale della città di Fabriano, con l'obiettivo di valutare il comportamento del modello in un contesto maggiormente rappresentativo di uno scenario applicativo concreto.

Un ulteriore obiettivo consiste nell'analizzare l'effetto del parametro $\gamma$, introdotto nel modello per rappresentare differenti politiche decisionali. La variazione di tale parametro consente di studiare il comportamento del sistema in presenza di una maggiore attenzione all'efficienza economica dell'azienda, di un compromesso neutrale tra i due obiettivi oppure di una maggiore priorità attribuita alla qualità del servizio e, conseguentemente, alla riduzione dell'insoddisfazione degli utenti.

Per ciascuna configurazione vengono analizzate sia la qualità della soluzione ottenuta sia le risorse necessarie per la sua realizzazione. In particolare, il confronto considera la funzione obiettivo complessiva e la sua scomposizione nelle componenti di insoddisfazione degli utenti e di costo logistico, oltre agli indicatori relativi alla saturazione fisica e temporale dei veicoli, al numero di mezzi impiegati e al tempo computazionale richiesto dagli algoritmi.

L'analisi sperimentale è quindi organizzata in una serie di benchmark, ciascuno finalizzato allo studio di uno specifico aspetto del problema. Tale struttura consente di valutare separatamente l'impatto della dimensione dell'istanza, della configurazione urbanistica, della distribuzione spaziale degli utenti e della differenza tra un grafo sintetico e una rete stradale reale. Due benchmark trasversali verificano infine su tutte le istanze l'effetto di $\gamma$ e il comportamento della soluzione a frequenza fissata. Il confronto sistematico tra le due euristiche permette infine di evidenziare eventuali differenze in termini di qualità delle soluzioni, efficienza nell'utilizzo della flotta e comportamento computazionale.


\section{Configurazione Sperimentale}\label{se:configurazione}



\subsection{Ambiente di Sviluppo}


La fase sperimentale è stata realizzata mediante una combinazione di strumenti software dedicati alla generazione delle istanze, alla risoluzione del problema e all'analisi dei risultati. La componente principale dell'implementazione è stata sviluppata in Python, utilizzando la versione 3.12.4.

Il modello e le euristiche di risoluzione sono state implementate in ambiente Python, facendo uso di diverse librerie per la gestione dei dati, il calcolo numerico, l'elaborazione dei grafi e la visualizzazione dei risultati. In particolare, sono state utilizzate:

\begin{itemize}
	\item \texttt{NumPy}, per le operazioni numeriche e la gestione vettorializzata delle matrici di distanza e di tempo;
	\item \texttt{SciPy}, per le procedure di calcolo scientifico e per le strutture e gli algoritmi utilizzati nella gestione delle reti;
	\item \texttt{Matplotlib}, per la produzione delle rappresentazioni grafiche dei risultati sperimentali;
	\item \texttt{OSMnx} e \texttt{GeoPandas}, per il preprocessing e la gestione dei dati geografici relativi alla rete stradale reale di Fabriano;
	\item \texttt{psutil}, utilizzata per il monitoraggio delle risorse hardware e per la determinazione dinamica del numero di worker impiegabili durante le elaborazioni.
\end{itemize}

L'analisi e la visualizzazione dei risultati dei benchmark sono state inoltre supportate da MATLAB, versione R2025a. Gli script MATLAB sono stati utilizzati per l'elaborazione dei risultati memorizzati nei file CSV e per la generazione dei grafici utilizzati nell'analisi sperimentale.

Le elaborazioni sono state eseguite su un sistema dotato di processore Apple M3, 24 GB di memoria RAM e sistema operativo macOS. La configurazione hardware viene riportata al fine di rendere esplicite le condizioni computazionali nelle quali sono stati misurati i tempi di esecuzione presentati nei benchmark.


\subsection{Parametri Utilizzati}



La configurazione sperimentale è stata definita mediante un insieme di parametri che caratterizzano i costi operativi, i vincoli temporali e di capacità, le frequenze di raccolta e le caratteristiche delle diverse tipologie di utenza. I valori utilizzati negli esperimenti sono riportati di seguito.

I principali parametri globali del modello sono riportati in Tabella~\ref{tab:parametri_globali}.

\newpage
\begin{table}[htbp]
	\centering
	\caption{Principali parametri globali utilizzati nelle simulazioni.}
	\label{tab:parametri_globali}
	\begin{tabular}{lll}
		\hline
		Parametro & Valore & Significato \\
		\hline
		$cd$ & $0.35$ \euro/km & Costo chilometrico \\
		$cm$ & $0.25$ \euro/min & Costo della manodopera \\
		$L$ & $480$ min & Durata massima del turno \\
		$\alpha$ & $10.0$ & Penalità per sottofrequenza \\
		$\beta$ & $2.0$ & Penalità per sovrafrequenza \\
		$X_{\mathrm{VALUES}}$ & $0.5,\ldots,6.0$ & Griglia delle frequenze \\
		$\gamma$ & $0.05,\ldots,0.95$ & Pesi (riferimento: $0.10$, $0.50$, $0.90$) \\
		\hline
	\end{tabular}
\end{table}

Le capacità dei veicoli e i relativi costi fissi di attivazione dipendono dalla tipologia di rifiuto considerata. I valori utilizzati sono riportati in Tabella~\ref{tab:veicoli_rifiuti}.

\begin{table}[htbp]
	\centering
	\caption{Capacità e costi fissi dei veicoli per tipologia di rifiuto.}
	\label{tab:veicoli_rifiuti}
	\begin{tabular}{lrr}
		\hline
		Tipologia di rifiuto & Capacità (kg) & Costo fisso (\euro) \\
		\hline
		Organico & 1500 & 120 \\
		Carta & 1500 & 80 \\
		Plastica & 1000 & 70 \\
		Vetro & 2000 & 110 \\
		Indifferenziata & 2000 & 90 \\
		\hline
	\end{tabular}
\end{table}


La quantità di rifiuto prodotta dagli utenti viene determinata a partire da un valore base associato alla tipologia di rifiuto e da un moltiplicatore dipendente dalla tipologia di utenza. I valori di riferimento sono riportati nelle Tabelle~\ref{tab:wbase} e~\ref{tab:moltiplicatori}.

\begin{table}[htbp]
	\centering
	\caption{Produzione base di rifiuto per tipologia.}
	\label{tab:wbase}
	\begin{tabular}{lr}
		\hline
		Tipologia di rifiuto & Produzione base (kg) \\
		\hline
		Organico & 3.0 \\
		Carta & 2.0 \\
		Plastica & 1.7 \\
		Vetro & 1.2 \\
		Indifferenziata & 2.5 \\
		\hline
	\end{tabular}
\end{table}

\begin{table}[htbp]
	\centering
	\caption{Moltiplicatori della produzione di rifiuto per tipologia di utenza.}
	\label{tab:moltiplicatori}
	\begin{tabular}{lr}
		\hline
		Tipologia di utenza & Moltiplicatore \\
		\hline
		Single & $0.5$ \\
		Famiglia & $1.0$ \\
		Palazzina piccola & $6.0$ \\
		Palazzina grande & $20.0$ \\
		\hline
	\end{tabular}
\end{table}

Allo stesso modo, la frequenza ideale $x^*_{rt}$ e il tempo di servizio per fermata $s_{rt}$ sono il prodotto di un valore base per rifiuto e di un moltiplicatore per tipologia di utenza (Tabella~\ref{tab:xstar_servizio}).

\begin{table}[htbp]
	\centering
	\caption{Frequenze ideali e tempi di servizio per fermata: valori base per rifiuto (sopra) e moltiplicatori per tipologia di utenza (sotto).}
	\label{tab:xstar_servizio}
	\begin{tabular}{lcc}
		\hline
		Rifiuto & $x^*$ base (ritiri/sett.) & $s$ base (min) \\
		\hline
		Organico        & 2.5 & 1.2 \\
		Carta           & 1.0 & 1.0 \\
		Plastica        & 1.7 & 1.0 \\
		Vetro           & 0.5 & 1.5 \\
		Indifferenziata & 2.0 & 1.2 \\
		\hline
	\end{tabular}
	
	\medskip
	\begin{tabular}{lcc}
		\hline
		Tipologia di utenza & Molt.\ $x^*$ & Molt.\ $s$ \\
		\hline
		Single            & 0.7 & 1.0 \\
		Famiglia          & 1.0 & 1.0 \\
		Palazzina piccola & 1.5 & 3.0 \\
		Palazzina grande  & 2.0 & 3.0 \\
		\hline
	\end{tabular}
\end{table}


Per rappresentare differenti configurazioni urbanistiche, sono state definite diverse distribuzioni percentuali delle tipologie di utenza. La composizione degli scenari utilizzati è riportata in Tabella~\ref{tab:scenari_urbanistici}.

\begin{table}[htbp]
	\centering
	\caption{Distribuzione delle tipologie di utenza nei diversi scenari urbanistici.}
	\label{tab:scenari_urbanistici}
	\begin{tabular}{lrrrr}
		\hline
		Scenario & Single & Famiglia & Pal. piccola & Pal. grande \\
		\hline
		Residenziale & 25\% & 45\% & 20\% & 10\% \\
		Suburbano & 20\% & 35\% & 30\% & 15\% \\
		Grandi condomini & 5\% & 10\% & 25\% & 60\% \\
		Villette & 40\% & 55\% & 5\% & 0\% \\
		Misto periferia & 15\% & 25\% & 45\% & 15\% \\
		\hline
	\end{tabular}
\end{table}


I parametri sopra riportati sono mantenuti costanti, salvo diversa indicazione, nelle diverse campagne sperimentali. La variazione dei parametri specificamente oggetto di ciascun benchmark viene descritta nelle relative sezioni, così da distinguere le condizioni comuni a tutti gli esperimenti dalle variabili utilizzate per analizzare i singoli fenomeni.






\subsection{Indicatori di prestazione}

Per valutare le soluzioni ottenute dalle due euristiche sono stati considerati diversi indicatori di prestazione, relativi sia alla qualità della soluzione sia all'utilizzo delle risorse e all'efficienza computazionale. Gli indicatori vengono calcolati durante l'esecuzione dell'algoritmo e successivamente aggregati negli script di analisi dei risultati.

\subsubsection{Funzione obiettivo}

La funzione obiettivo complessiva $F_{\mathrm{total}}$ rappresenta il principale indicatore di qualità della soluzione. Essa combina la componente relativa all'insoddisfazione degli utenti con la componente relativa ai costi logistici:

\[
F_{\mathrm{total}}
=
\gamma \left(F_{\mathrm{insoddis}}\, k_{\mathrm{scala}}\right)
+
(1-\gamma)F_{\mathrm{costi}},
\]

dove $F_{\mathrm{insoddis}}$ rappresenta l'insoddisfazione complessiva, $F_{\mathrm{costi}}$ rappresenta i costi logistici complessivi e $k_{\mathrm{scala}}$ è un fattore di bilanciamento utilizzato per rendere confrontabili le due componenti della funzione obiettivo.

La componente dei costi logistici è a sua volta costituita dai costi fissi di attivazione dei veicoli, dai costi di viaggio e dai costi della manodopera:

\[
F_{\mathrm{costi}}
=
F_{\mathrm{costo\_fisso}}
+
F_{\mathrm{viaggio}}
+
F_{\mathrm{lavoro}}.
\]

Il valore di $F_{\mathrm{total}}$ viene utilizzato dalla procedura di ricerca per individuare, tra i valori della griglia di frequenze considerati, la configurazione che minimizza la funzione obiettivo.

\subsubsection{Insoddisfazione degli utenti}

L'insoddisfazione complessiva, indicata con $F_{\mathrm{insoddis}}$, misura la deviazione tra la frequenza di raccolta programmata $X_r$ e la frequenza ideale $x^*_{ru}$ associata a ciascun utente. Nel modello la penalizzazione è asimmetrica e distingue tra sottofrequenza e sovrafrequenza:

\[
F_{\mathrm{insoddis}}
=
\sum_u
\left[
\alpha \max(0,x^*_{ru}-X_r)
+
\beta \max(0,X_r-x^*_{ru})
\right].
\]

Con i valori adottati nella sperimentazione, $\alpha=10$ e $\beta=2$, una frequenza inferiore a quella ideale viene quindi penalizzata maggiormente rispetto a una frequenza superiore. Tale scelta riflette l'ipotesi secondo cui il sotto-servizio risulta più penalizzante rispetto all'effettuazione di ritiri aggiuntivi.

\subsubsection{Costi logistici}

I costi logistici comprendono le tre componenti principali considerate dal modello: costo fisso di attivazione dei veicoli, costo di percorrenza e costo della manodopera.

Il costo fisso dipende dal numero di veicoli attivati e dal costo specifico associato alla tipologia di rifiuto. Il costo di viaggio dipende invece dalla distanza complessivamente percorsa dai veicoli, mentre il costo della manodopera è determinato dal tempo complessivamente impiegato e dal relativo costo per unità di tempo.

L'indicatore $F_{\mathrm{costi}}$ è quindi utilizzato per quantificare l'impegno economico necessario per realizzare il servizio di raccolta associato alla soluzione considerata.

\subsubsection{Saturazione fisica}

La saturazione fisica misura il grado di utilizzo della capacità di carico della flotta. È calcolata come rapporto tra il carico complessivamente trasportato e la capacità teorica complessiva dei veicoli attivati:

\[
\mathrm{Sat}_{\mathrm{fisica}}
=
\frac{\sum \mathrm{load}}
{n_{\mathrm{vehicles}}\,C_r}
\cdot 100.
\]

Il valore è espresso in percentuale. Un valore elevato indica un maggiore sfruttamento della capacità disponibile, mentre valori più bassi indicano una capacità residua maggiore.

\subsubsection{Saturazione temporale}

La saturazione temporale misura il grado di utilizzo del tempo massimo disponibile per la flotta. Essa è definita come:

\[
\mathrm{Sat}_{\mathrm{tempo}}
=
\frac{\sum \mathrm{time}}
{n_{\mathrm{vehicles}}\,L}
\cdot 100,
\]

dove $L$ rappresenta la durata massima del turno di lavoro.

Anche questo indicatore è espresso in percentuale e consente di valutare quanto la soluzione sfrutti il tempo disponibile per l'esecuzione dei percorsi.

\subsubsection{Numero di veicoli}

Il numero di veicoli $V_r$ è il numero di camion necessari per un giro di raccolta del rifiuto $r$; moltiplicato per la frequenza $X_r$ dà le attivazioni settimanali, che determinano i costi fissi di attivazione. L'indicatore permette di valutare l'impatto delle diverse configurazioni del problema sull'impiego delle risorse aziendali.

\subsubsection{Tempo computazionale}

Il tempo computazionale misura il tempo di elaborazione necessario all'algoritmo per completare la procedura di ricerca della soluzione. Nel file dei risultati tale valore è indicato dal campo \texttt{algo\_time\_sec} ed è espresso in secondi.

Il tempo viene misurato mediante un riferimento temporale di tipo wall-clock e comprende l'intera procedura di ricerca effettuata dall'algoritmo sull'istanza considerata, compreso l'avvio dei processi paralleli. Tale indicatore è particolarmente rilevante nel confronto tra Greedy e Clarke-Wright e nell'analisi della scalabilità rispetto alla dimensione dell'istanza.

\subsubsection{Soluzione selezionata}

Per distinguere la soluzione selezionata dalle altre configurazioni esplorate durante la ricerca, i risultati riportano inoltre il campo \texttt{is\_best}. Esso assume valore pari a $1$ per la configurazione di frequenza $X_r$ che minimizza $F_{\mathrm{total}}$ all'interno della griglia considerata, mentre assume valore pari a $0$ per le altre configurazioni.

Questo indicatore consente di identificare nei file CSV la soluzione utilizzata per l'analisi dei benchmark e di ricostruire, quando necessario, l'andamento della funzione obiettivo rispetto ai valori di frequenza esplorati.

\subsection{Metodo di analisi dei risultati}\label{se:metodo-analisi}

\paragraph{Ruolo di $\gamma$.}
Per ogni valore fissato di $X_r$ le euristiche di routing minimizzano il costo logistico, mentre la componente di insoddisfazione $F_{\mathrm{insoddis}}(X_r)$ dipende soltanto da $X_r$ e non dalle rotte. A parità di $X_r$, quindi, le rotte costruite, il numero di veicoli e tutte le componenti non pesate della soluzione non dipendono da $\gamma$. Il parametro interviene esclusivamente nella Grid Search, dove seleziona
\begin{equation}
	X_r^{*}(\gamma)=\arg\min_{X_r\in\mathcal{X}}\;\Bigl[\gamma\,k_{\mathrm{scala}}\,F_{\mathrm{insoddis}}(X_r)+(1-\gamma)\,F_{\mathrm{costi}}(X_r)\Bigr].
\end{equation}
La proprietà è stata verificata numericamente. Per tutte le 1920 combinazioni (istanza, algoritmo, rifiuto, $X_r$) dei benchmark originali rotte e componenti coincidono tra i tre valori di $\gamma$. Ricostruendo $X_r^*$ a posteriori dalle componenti grezze si ritrovano le soluzioni selezionate in 480 casi su 480, e rieseguendo l'intera procedura con $\gamma=0.3$ e $\gamma=0.7$ (valori mai usati nelle run originali) le frequenze ottime coincidono con quelle previste in 20 casi su 20. Per questo i risultati per un qualunque valore di $\gamma$ si ottengono dalle stesse esecuzioni, senza ripetere il routing. I benchmark usano come riferimento le tre politiche $\gamma=0.10$, $0.50$ e $0.90$; le analisi in funzione di $\gamma$ (Benchmark~1 e~4 e benchmark trasversali) usano i 19 valori $\gamma=0.05,0.10,\dots,0.95$.

Ne segue che il confronto tra i valori di $F_{\mathrm{tot}}$ al variare di $\gamma$ non è informativo: $F_{\mathrm{tot}}$ cambia anche solo perché cambiano i pesi, e il fattore $k_{\mathrm{scala}}$ varia da istanza a istanza (da $0.55$ a $2.05$ tra le istanze dei benchmark originali). Le politiche vengono quindi confrontate sulla soluzione fisica che producono.

\paragraph{Grandezze fisiche.}
Per ogni soluzione si considerano, sommate sui cinque rifiuti e non pesate:
{\sloppy
	\begin{itemize}
		\item il numero di camion per giro $V=\sum_r V_r$ e le attivazioni settimanali $\sum_r V_r X_r$ (il modello non include un calendario settimanale: $V_r$ è la dimensione della flotta di un singolo giro, non il parco veicoli);
		\item i chilometri settimanali $\sum_r F_{\mathrm{viaggio}}/cd$ e le ore di lavoro settimanali $\sum_r F_{\mathrm{lavoro}}/(60\,cm)$;
		\item il disservizio grezzo $F_{\mathrm{insoddis}}$, non scalato, e il costo logistico grezzo $F_{\mathrm{costi}}$ in euro;
		\item le frequenze ottime $X_r^*$ e le saturazioni fisica e temporale.
	\end{itemize}
	\par}
Il tempo di servizio per fermata dipende dal rifiuto e dalla tipologia di utenza (Tabella~\ref{tab:xstar_servizio}) e non da $X_r$; il tempo di viaggio è proporzionale ai chilometri (velocità di 25\,km/h).

\paragraph{Gap tra Greedy e Clarke-Wright.}
Il gap è calcolato come
\begin{equation}
	\text{Gap}\,\% = \frac{F^{G}_{\mathrm{tot}}-F^{CW}_{\mathrm{tot}}}{\min\!\bigl(F^{G}_{\mathrm{tot}},F^{CW}_{\mathrm{tot}}\bigr)}\cdot 100,
\end{equation}
sulle soluzioni selezionate dai due algoritmi allo stesso $\gamma$. Poiché i due algoritmi possono scegliere frequenze diverse, la differenza di $F_{\mathrm{tot}}$ si scompone in
\begin{equation}
	F^{G}(X^{*G})-F^{CW}(X^{*CW}) = \underbrace{\bigl[F^{G}(X^{*G})-F^{CW}(X^{*G})\bigr]}_{\text{routing a pari frequenza}} + \underbrace{\bigl[F^{CW}(X^{*G})-F^{CW}(X^{*CW})\bigr]}_{\text{spostamento di }X^*}.
\end{equation}
A parità di frequenze il gap risulta inoltre diluito dal termine di insoddisfazione, che è identico per i due algoritmi:
\begin{equation}
	\text{Gap}\,\% \simeq \frac{(1-\gamma)\,\Delta F_{\mathrm{costi}}}{F_{\mathrm{tot}}}\cdot 100 .
\end{equation}
Per questo motivo il gap su $F_{\mathrm{tot}}$ non è una misura diretta del vantaggio di routing: l'analisi affianca sempre il gap logistico a $X_r$ fissato.

\paragraph{Replicazione.}
Le istanze sintetiche sono replicate su 30 seed estratti casualmente e si riportano medie e intervalli di confidenza al 95\% (distribuzione $t$ Student). La mappa di Fabriano è una singola istanza deterministica, che include tutti gli utenti reali, e non ha intervallo di confidenza; la variabilità dovuta alla scelta degli utenti è stimata con 30 sottocampioni casuali di utenti reali (Sezione~\ref{sse:b4-sintetico-reale}).


\section{Benchmark 1 - Analisi della Variazione del Numero di Utenti}\label{se:b1}

\subsection{Descrizione dello scenario}

Il primo benchmark analizza come cambiano la soluzione e il vantaggio di Clarke-Wright rispetto a Greedy al crescere della dimensione dell'istanza, espressa dal numero di utenti attivi
\[
N \in \{100,250,500,750,1000\}.
\]
Le istanze sono sintetiche: per ciascun seed viene generata una rete su un'area di $10\times10$\,km con il deposito al centro, da cui vengono estratti a caso gli $N$ utenti attivi (distribuzione uniforme, scenario urbanistico \textit{residenziale}). A parità di $N$ e di seed, i due algoritmi lavorano sulla stessa istanza, per cui le differenze sono attribuibili solo alla procedura di costruzione delle rotte.

Come descritto nella Sezione~\ref{se:metodo-analisi}, per ogni istanza ed ogni algoritmo si risolve una sola griglia di frequenze, e le tre politiche $\gamma\in\{0.10,\,0.50,\,0.90\}$ (Pro-Azienda, Neutra, Pro-Cittadino) vengono applicate in post-elaborazione scegliendo la frequenza che minimizza $F_{\mathrm{tot}}$. Poiché $k_{\mathrm{scala}}$ viene ricalcolato per ogni istanza (da $0.84$ a $2.05$ in questo benchmark), i valori assoluti di $F_{\mathrm{tot}}$ non sono confrontabili tra dimensioni diverse e non vengono usati: l'analisi si basa sulla soluzione fisica (camion per giro, attivazioni, chilometri, ore di lavoro, saturazioni, disservizio non pesato), sul gap tra i due algoritmi e sul tempo di calcolo. Le grandezze sono medie su 30 seed con intervallo di confidenza al $95\%$.

\subsection{Discussione dei risultati}\label{sse:b1-discussione}

In tutte le istanze del benchmark tutti gli utenti sono serviti ($5N$ servizi per giro completo): la copertura non distingue né gli algoritmi né le dimensioni, e il confronto si concentra sulle grandezze fisiche e sul gap.

\begin{table}[htbp]
	\centering
	\caption{Benchmark 1: gap su $F_{\mathrm{tot}}$ (media su 30 seed) per politica; variazione del numero di camion di CW rispetto a Greedy a $\gamma=0.5$; vantaggio logistico di CW a pari $X_r$ nel regime a tempo ($X_r\ge3$), con ripartizione approssimata tra chilometri e camion (punti percentuali).}
	\label{tab:b1_gap}
	\small
	\resizebox{\textwidth}{!}{%
		\begin{tabular}{rcccrrrr}
			\hline
			$N$ & Gap $\gamma=0.1$ & Gap $\gamma=0.5$ & Gap $\gamma=0.9$ & $\Delta V$ (\%) & Gap logistico & di cui km & di cui camion \\
			\hline
			100  & $11.81\pm0.99$ & $21.67\pm2.36$ & $5.99\pm0.57$ & $-35.0$ & 42.8 & 19.2 & 23.6 \\
			250  & $7.45\pm0.57$  & $13.82\pm1.53$ & $4.20\pm0.44$ & $-18.3$ & 27.4 & 18.1 & 9.3 \\
			500  & $7.16\pm0.52$  & $8.70\pm0.54$  & $2.77\pm0.17$ & $-6.5$  & 19.2 & 16.7 & 2.5 \\
			750  & $7.17\pm0.35$  & $10.65\pm0.62$ & $3.48\pm0.18$ & $-10.2$ & 21.2 & 16.0 & 5.2 \\
			1000 & $7.46\pm0.32$  & $9.42\pm0.47$  & $3.14\pm0.16$ & $-7.0$  & 19.1 & 15.7 & 3.4 \\
			\hline
		\end{tabular}%
	}
\end{table}

\paragraph{Andamento del gap rispetto a $N$.}\mbox{}\\
A $\gamma=0.5$ il gap passa da $21.67\pm2.36\%$ a $N=100$ a $13.82\pm1.53\%$ a $N=250$ e a $8.70\pm0.54\%$ a $N=500$, risale a $10.65\pm0.62\%$ a $N=750$ e scende a $9.42\pm0.47\%$ a $N=1000$ (Tabella~\ref{tab:b1_gap}). Il minimo a $N=500$ e la risalita a $750$ sono differenze reali, oltre quattro volte l'errore standard, e non rumore di campionamento; tra $N=500$ e $N=1000$, tuttavia, il gap resta in una fascia stretta ($8.7$--$10.7\%$), molto inferiore a quella di $N=100$.

Il vantaggio logistico a pari frequenza (regime a tempo) si scompone in un termine di chilometri e in uno di camion ($\Delta V\,c_r\,X_r$). Il termine di chilometri varia poco ($19.2$, $18.1$, $16.7$, $16.0$ e $15.7$ punti): il risparmio relativo di CW cresce con $N$ (il rapporto Greedy/CW sui chilometri passa da $1.43$ a $1.74$, cioè da $-30\%$ a $-42\%$ di percorrenza), ma il peso del viaggio cala (dal $59.8\%$ al $28.0\%$ del tempo di giro di CW) e i due effetti si compensano quasi. Il termine di camion decade con $N$ ($23.6$, $9.3$, $2.5$, $5.2$, $3.4$ punti) perché lo stesso scarto pesa meno su una flotta più grande, ed è non monotono perché $\Delta V$ è un numero intero ($3.5$, $2.7$, $1.4$, $2.9$ e $2.4$ camion in media): la forma del gap segue la riduzione percentuale di camion di CW ($-35.0\%$, $-18.3\%$, $-6.5\%$, $-10.2\%$, $-7.0\%$), minima a $N=500$ e massima a $N=750$ tra le istanze maggiori.

\begin{figure}[htbp]
	\centering
	\includegraphics[width=0.75\textwidth]{immagini/fig_gap_vs_N_B1.png}
	\caption{Benchmark 1: gap su $F_{\mathrm{tot}}$ tra Greedy e Clarke-Wright al variare di $N$ per le tre politiche (media su 30 seed).}
	\label{fig:b1_gap}
\end{figure}

\paragraph{Effetto di $\gamma$ sul gap.}\mbox{}\\
A $\gamma=0.9$ i due algoritmi scelgono le stesse frequenze e il gap è routing puro diluito dal termine di insoddisfazione ($5.99\%$ a $N=100$ e $2.8$--$4.2\%$ per $N\ge250$). A $\gamma=0.1$ il gap per $N\ge250$ non dipende da $N$ ($7.2$--$7.5\%$, intervalli sovrapposti) e una quota del $12$--$30\%$ è dovuta allo spostamento di $X^*$ ($2.25$ punti a $N=250$, $0.87$ a $N=1000$). Il picco a $\gamma=0.5$ vale per tutte le dimensioni, con intervalli non sovrapposti rispetto a $\gamma=0.1$ e $0.9$: a $\gamma=0.5$ la frequenza ottima sta sopra la soglia tra i due regimi (Sezione~\ref{sse:t2-regimi}), dove il gap logistico è maggiore (circa $13\%$ a $X_r=0.5$ per $N\ge250$, contro $19$--$27\%$ a $X_r\ge3$), e il peso del costo $1-\gamma$ è ancora $0.5$.


\begin{figure}[htbp]
	\centering
	\includegraphics[width=0.7\textwidth]{immagini/fig_gap_profilo_B1.png}
	\caption{Benchmark 1: gap logistico tra Greedy e CW (fisso, viaggio e lavoro) a frequenza $X_r$ comune, calcolato per rifiuto e mediato sui rifiuti e sui 30 seed, per quattro valori di $N$ (fasce IC95).}
	\label{fig:b1_profilo}
\end{figure}

Il gap logistico è basso sotto capacità e raggiunge un plateau nel regime a tempo ($42.8\%$ a $N=100$, tra $19.1$ e $27.4\%$ per $N\ge250$, Tabella~\ref{tab:b1_gap}). Il picco del gap su $F_{\mathrm{tot}}$ a $\gamma=0.5$ nasce perché $X^*$ cade sopra questa soglia, dove il vantaggio di CW è maggiore.

\paragraph{Il paradosso dei camion.}\mbox{}\\
A $N=1000$ il numero di camion per giro a $\gamma=0.1$ è $49.0$ per Greedy e $38.9$ per CW, contro $34.1$ e $31.7$ a $\gamma=0.5$; le attivazioni settimanali vanno in senso opposto ($29.0$ contro $56.6$ per Greedy, $30.2$ contro $58.1$ per CW). La saturazione fisica è $86.8\%$ (Greedy) e $83.3\%$ (CW) a $\gamma=0.1$ e $50.2\%$ e $48.8\%$ a $\gamma=0.5$, la temporale $72.9\%$ e $73.7\%$ contro $95.7\%$ e $85.6\%$: a frequenze basse il carico per ritiro $Q_{ru}=W_{rt_u}/X_r$ rende vincolante la capacità, a frequenze alte il tempo del turno. Il fenomeno cresce con la dimensione: il rapporto tra camion a $\gamma=0.1$ e $0.5$ per Greedy vale $1.00$, $1.15$, $1.33$, $1.29$ e $1.44$ per $N=100,\dots,1000$; a $N=100$ non c'è, perché la capacità non è mai vincolante (V resta $10$ per Greedy e $6.5$ per CW a ogni $\gamma$). Tra $\gamma=0.5$ e $\gamma=0.9$ il numero di camion è identico per entrambi gli algoritmi e per ogni $N$ (a $N=1000$ $34.1$ e $31.7$), mentre il lavoro cresce: le attivazioni settimanali passano da $56.6$ a $80.4$ per Greedy e da $58.1$ a $74.8$ per CW, i chilometri da $4380$ a $6172$ e da $2801$ a $3569$, le ore da $431$ a $613$ e da $395$ a $510$. La flotta è resiliente al peso, il lavoro no.

\begin{figure}[htbp]
	\centering
	\includegraphics[width=\textwidth]{immagini/fig_fisica_B1.png}
	\caption{Benchmark 1 ($N=1000$): soluzione selezionata al variare di $\gamma$ (media su 30 seed, fasce IC95): camion per giro, attivazioni settimanali, chilometri, ore di lavoro, disservizio grezzo e frequenze ottime $X^*_r$ per rifiuto (CW continuo, Greedy tratteggiato).}
	\label{fig:b1_fisica}
\end{figure}

Il numero di camion crolla tra $\gamma=0.05$ e $0.15$ e poi resta costante ($34.1$ per Greedy e $31.7$ per CW da $\gamma=0.5$), mentre attivazioni, chilometri e ore continuano a salire a gradini: la flotta è resiliente al peso, il lavoro no. I gradini coincidono con i cambi delle frequenze ottime (ultimo pannello).

\paragraph{Costo marginale e scelta di $\gamma$.}\mbox{}\\
Il costo marginale per unità di disservizio eliminata tra $\gamma=0.5$ e $0.9$ è $6.9$--$8.2$ volte quello tra $0.1$ e $0.5$ per CW ($6.5$--$7.4$ per Greedy), senza dipendenza evidente da $N$: a $N=1000$, per CW, da $0.196$\,\euro\ a $1.358$\,\euro\ per unità. Con tre soli valori di $\gamma$ il salto sembra cadere a $\gamma=0.5$; sulla griglia di 19 valori il costo marginale cresce invece a ogni passo e il salto più grande cade tra $\gamma=0.85$ e $0.90$ a tutte le dimensioni (Sezione~\ref{sse:t1-soglie}).

\paragraph{Clarke-Wright sotto capacità.}\mbox{}\\
A $X_r=0.5$, dove vincola la capacità, CW non riduce il numero di camion: per $N=250,\dots,1000$ ne usa in media $0.6$--$0.9$ in più di Greedy (differenza Greedy$-$CW pari a $-0.77\pm0.35$, $-0.73\pm0.49$, $-0.90\pm0.43$, $-0.57\pm0.47$), mentre a $N=100$ ne usa $1.17\pm0.44$ in meno. Coerentemente, l'eccesso rispetto al limite inferiore di capacità $\lceil\sum Q/C\rceil$, sommato sui rifiuti, è maggiore per CW da $N=250$ in su ($1.87$ contro $1.10$, $2.77$ contro $2.03$, $3.70$ contro $2.80$, $3.60$ contro $3.03$ camion) e minore a $N=100$ ($1.27$ contro $2.43$). Il meccanismo non è stato isolato: si registra la differenza, non la sua causa.

\paragraph{Tempi di calcolo.}\mbox{}\\
La Tabella~\ref{tab:b1_tempi} riporta il tempo di una griglia completa ($12$ valori di $X_r$ per $5$ rifiuti, con processi paralleli sulla macchina descritta nella Sezione~\ref{se:metodo-analisi}), medio su 30 seed. Fino a $N=250$ i due algoritmi richiedono lo stesso tempo, circa $9.5$\,s: è una componente fissa, indipendente da $N$ e dall'algoritmo (il tempo a $N=100$ è $9.57$\,s per Greedy e $9.58$\,s per CW), che comprende l'avvio dei processi paralleli, ripetuto per ogni rifiuto: avviandoli una sola volta per l'intera griglia, il tempo a $N=100$ scende a circa $5.4$\,s (Sezione~\ref{sse:t2-potatura}). Oltre $N=250$ i tempi crescono: sopra la componente fissa l'incremento cresce circa come $N^2$ per entrambi (esponente $2.0$ per CW e $1.95$ per Greedy, stimato su $N\ge500$), con una costante circa tre volte maggiore per CW (rapporto tra incrementi $2.6$--$3.6$); poiché l'incremento comprende anche il trasferimento dei dati ai processi, uguale per i due algoritmi, questo rapporto sottostima la differenza tra i soli tempi di routing. Il rapporto tra i tempi totali passa così da $1.00$ a $N=100$ a $1.37$ a $N=1000$ e a $2.16$ a $N=2000$ (campionamenti di Fabriano); sull'istanza completa di Fabriano ($N=3932$, campagna principale) è $2.38$ ($100.8$\,s contro $42.3$\,s), mentre misurando i due algoritmi con un unico insieme di processi è $5.2$ ($87.6$\,s contro $16.8$\,s, Sezione~\ref{sse:t2-potatura}). I tempi non dipendono dalla natura dell'istanza: a $N=1000$ valgono $15.73$\,s (CW) e $11.49$\,s (Greedy) sul sintetico e $15.63$\,s e $11.34$\,s sui sottocampioni reali, e a $N=400$ sono $10.4$--$10.6$\,s e $9.9$--$10.0$\,s in tutte le tipologie e distribuzioni. In valore assoluto entrambi sono praticamente utilizzabili: una griglia completa su $3932$ utenti richiede meno di due minuti.

\begin{table}[htbp]
	\centering
	\caption{Tempo di una griglia completa (secondi, media su 30 seed) e rapporto CW/Greedy. $^{*}$: campagna principale, singola esecuzione.}
	\label{tab:b1_tempi}
	\small
	\begin{tabular}{lrrrr}
		\hline
		Istanza & $N$ & Greedy & CW & CW/Greedy \\
		\hline
		Sintetico (B1)      & 100  & 9.57  & 9.58  & 1.00 \\
		Sintetico (B1)      & 250  & 9.71  & 9.83  & 1.01 \\
		Sintetico (B1)      & 500  & 10.14 & 11.06 & 1.09 \\
		Sintetico (B1)      & 750  & 10.82 & 12.94 & 1.20 \\
		Sintetico (B1)      & 1000 & 11.49 & 15.73 & 1.37 \\
		Fabriano campionato & 2000 & 16.89 & 36.46 & 2.16 \\
		Fabriano completo$^{*}$ & 3932 & 42.3 & 100.8 & 2.38 \\
		\hline
	\end{tabular}
\end{table}

\paragraph{Sintesi.}\mbox{}\\
Il benchmark mostra tre fatti: la dimensione dell'istanza sposta il gap soprattutto attraverso il termine discreto dei camion e il peso del viaggio, non attraverso la qualità relativa dei chilometri (che migliora con $N$); $\gamma$ cambia la soluzione fisica in modo costante a tratti, con un costo marginale che cresce a ogni passo; e il vantaggio di CW sui camion esiste nel regime a tempo e non sotto capacità.



\section{Benchmark 2 - Analisi delle tipologie urbanistiche}\label{se:b2}



\subsection{Descrizione degli scenari}

Il secondo benchmark è finalizzato ad analizzare l'influenza della tipologia urbanistica sulla qualità delle soluzioni ottenute per il problema di raccolta differenziata. A differenza del Benchmark 1, in cui viene studiato l'effetto della variazione del numero di utenti, in questa analisi il numero di utenti viene mantenuto invariato e viene modificata la composizione della popolazione in funzione della struttura insediativa considerata.

Sono stati considerati cinque differenti scenari urbanistici: \textit{residenziale}, \textit{suburbano}, \textit{grandi condomini}, \textit{villette} e \textit{misto periferia}. Ciascuno scenario è caratterizzato da una diversa distribuzione delle quattro tipologie di utenza introdotte nel modello, ovvero \textit{Single}, \textit{Famiglia}, \textit{Palazzina piccola} e \textit{Palazzina grande}. La variazione di tali proporzioni permette di rappresentare configurazioni urbane caratterizzate da differenti livelli di concentrazione degli utenti e, di conseguenza, da differenti quantità di rifiuti prodotti e modalità di organizzazione della raccolta.

In particolare, lo scenario \textit{residenziale} presenta una distribuzione maggiormente equilibrata tra famiglie e altre tipologie di utenza, mentre lo scenario \textit{suburbano} è caratterizzato da una maggiore presenza di palazzine. Lo scenario \textit{grandi condomini} rappresenta invece una configurazione fortemente concentrata, nella quale prevalgono le palazzine di grandi dimensioni. Al contrario, lo scenario \textit{villette} descrive una realtà a bassa densità abitativa, caratterizzata prevalentemente da utenti singoli e famiglie. Infine, lo scenario \textit{misto periferia} rappresenta una situazione intermedia, con una presenza significativa di palazzine piccole e una distribuzione più eterogenea delle utenze.

Le percentuali utilizzate per costruire i cinque scenari sono riportate nella Tabella~\ref{tab:scenari_urbanistici} (Sezione~\ref{se:configurazione}).

Per ciascuna configurazione urbanistica, su $N=400$ utenti distribuiti uniformemente, vengono applicati entrambi gli algoritmi su 30 seed. Come nel Benchmark 1, le tre politiche $\gamma=0.10$ (Pro-Azienda), $0.50$ (Neutra) e $0.90$ (Pro-Cittadino) sono applicate in post-elaborazione sulla stessa griglia di frequenze (Sezione~\ref{se:metodo-analisi}). Poiché $k_{\mathrm{scala}}$ dipende dalla composizione degli utenti, i valori di $F_{\mathrm{tot}}$ non sono confrontabili tra scenari: il confronto tra tipologie si basa sulla soluzione fisica (camion, attivazioni, costo logistico e disservizio non pesati, saturazioni, frequenze ottime) e sul gap tra i due algoritmi.

\newpage
\subsection{Discussione dei risultati}\label{sse:b2-discussione}

\begin{table}[htbp]
	\centering
	\caption{Benchmark 2 ($N=400$): gap su $F_{\mathrm{tot}}$ (media $\pm$ IC95 su 30 seed) per scenario e politica.}
	\label{tab:b2_gap}
	\small
	\begin{tabular}{lccc}
		\hline
		Scenario & $\gamma=0.1$ & $\gamma=0.5$ & $\gamma=0.9$ \\
		\hline
		Residenziale     & $6.17\pm0.48$ & $11.58\pm1.04$ & $3.69\pm0.28$ \\
		Suburbano        & $6.88\pm0.63$ & $11.97\pm1.01$ & $3.45\pm0.24$ \\
		Grandi condomini & $4.03\pm0.28$ & $12.38\pm1.50$ & $6.40\pm0.74$ \\
		Villette         & $7.31\pm0.50$ & $12.10\pm0.80$ & $4.06\pm0.28$ \\
		Misto periferia  & $5.82\pm0.36$ & $9.68\pm0.73$  & $2.99\pm0.22$ \\
		\hline
	\end{tabular}
\end{table}


\begin{figure}[htbp]
	\centering
	\includegraphics[width=0.9\textwidth]{immagini/fig_gap_B2.png}
	\caption{Benchmark 2: gap su $F_{\mathrm{tot}}$ tra Greedy e Clarke-Wright per scenario urbanistico e politica (media su 30 seed, barre: IC95).}
	\label{fig:b2_gap}
\end{figure}


\begin{table}[htbp]
	\centering
	\caption{Benchmark 2: soluzione di CW per scenario (media su 30 seed). $V$ a $\gamma=0.1\rightarrow0.5$; attivazioni, costo logistico grezzo e disservizio grezzo a $\gamma=0.5$; saturazione fisica a $\gamma=0.1$; frequenze ottime di indifferenziata e organico a $\gamma=0.5$.}
	\label{tab:b2_fisica}
	\small
	\begin{tabular}{lcrrrrcc}
		\hline
		Scenario & $V$ & $V\!\cdot\!X$ & $F_{\mathrm{costi}}$ (\euro) & $F_{\mathrm{ins}}$ & Sat.\ fis.\ (\%) & $X^*_{ind}$ & $X^*_{org}$ \\
		\hline
		Residenziale     & $21.6\rightarrow16.5$ & 30.9 & 6568  & 5250 & 76.6 & 2.6 & 2.8 \\
		Suburbano        & $24.8\rightarrow18.2$ & 39.7 & 8367  & 5114 & 82.8 & 3.0 & 3.6 \\
		Grandi condomini & $69.7\rightarrow22.7$ & 69.2 & 14510 & 1858 & 94.1 & 4.0 & 5.0 \\
		Villette         & $14.2\rightarrow14.2$ & 22.1 & 4613  & 1772 & 39.2 & 2.0 & 2.5 \\
		Misto periferia  & $29.6\rightarrow19.9$ & 47.3 & 9998  & 4030 & 84.7 & 3.0 & 4.0 \\
		\hline
	\end{tabular}
\end{table}

\paragraph{La tipologia cambia la soluzione fisica più del gap.}\mbox{}\\
Il carico per fermata e le frequenze ideali determinano la soluzione (Tabella~\ref{tab:b2_fisica}). A $\gamma=0.5$ il costo logistico grezzo di CW è $14\,510$\,\euro\ nei grandi condomini contro $4613$\,\euro\ nelle villette ($3.1$ volte): le frequenze ottime seguono le frequenze ideali (indifferenziata a $4.0$ e organico a $5.0$ ritiri nei grandi condomini, a $2.0$ e $2.5$ nelle villette) e il carico per ritiro rende la capacità vincolante a $\gamma=0.1$ (saturazione fisica $94.1\%$ contro $39.2\%$). Il gap su $F_{\mathrm{tot}}$ a $\gamma=0.5$, invece, è praticamente lo stesso in quattro scenari su cinque ($11.58$--$12.38\%$, intervalli sovrapposti) e solo il misto è più basso ($9.68\pm0.73\%$): l'ordinamento degli scenari osservato con un solo seed (grandi condomini massimo con $16.15\%$) era un effetto dell'istanza e non si ripete con 30 seed ($12.38\pm1.50\%$).

\paragraph{Il paradosso dei camion dipende dal carico per fermata.}\mbox{}\\
Nei grandi condomini il numero di camion di Greedy a $\gamma=0.1$ è $77.9$ contro $23.6$ a $\gamma=0.5$ (CW $69.7$ contro $22.7$), un rapporto di $3.3$; le attivazioni settimanali vanno da $38.9$ a $72.7$ (Greedy). Negli altri scenari il rapporto è $1.3$--$1.7$ e nelle villette non esiste: $V$ resta $15.0$ (Greedy) e $14.2$ (CW) a ogni $\gamma$ perché la capacità non è mai vincolante (saturazione fisica $14$--$39\%$). La stessa meccanica $Q_{ru}=W_{rt_u}/X_r$ spiega i due estremi: più carico per fermata, più camion necessari quando la frequenza scende.

\begin{figure}[htbp]
	\centering
	\includegraphics[width=\textwidth]{immagini/fig_V_vs_X.png}
	\caption{Camion per giro dell'organico in funzione della frequenza $X_r$ (media su 30 seed, fasce IC95): $B1$ con $N=1000$, grandi condomini e villette con $N=400$. La linea tratteggiata è $V\cdot X$ costante, ancorata al valore di Greedy a $X_r=0.5$.}
	\label{fig:b2_V_vs_X}
\end{figure}

La Figura~\ref{fig:b2_V_vs_X} mostra la meccanica $Q_{ru}=W_{rt_u}/X_r$: sotto la soglia di capacità il numero di camion segue l'iperbole $V\cdot X\approx$ costante, sopra la soglia diventa indipendente da $X_r$ perché vincola il tempo. Più alto è il carico per fermata (grandi condomini), più alta è la soglia e più è ripido il crollo dei camion quando la frequenza sale; nelle villette la capacità non vincola.


\paragraph{Resilienza e disservizio residuo.}\mbox{}\\
Nei grandi condomini la soluzione a $\gamma=0.5$ coincide con quella a $\gamma=0.9$ (stesse frequenze: $2.0$, $4.0$, $5.0$, $3.5$, $1.0$): il disservizio grezzo vale $1858$ a entrambi i valori e non può scendere oltre. Poiché $X_r$ è comune a tutti gli utenti, il disservizio non si annulla quando le frequenze ideali sono eterogenee; il valore residuo a $\gamma=0.5$ è più basso negli scenari omogenei (villette $1772$, grandi condomini $1858$) e più alto in quelli misti ($4030$--$5250$). Per lo stesso motivo il costo marginale del secondo tratto ($0.5\rightarrow0.9$) non è definito nei grandi condomini; negli altri scenari è $5.7$--$7.6$ volte quello del primo tratto per CW (misto $10.3$).

\paragraph{Gap nei grandi condomini a $\gamma=0.1$ e $0.9$.}\mbox{}\\
I grandi condomini sono l'unico scenario in cui il gap a $\gamma=0.1$ ($4.03\pm0.28\%$) è molto più basso di quello a $\gamma=0.9$ ($6.40\pm0.74\%$). A $\gamma=0.1$ tutte le frequenze stanno a $0.5$ (capacità vincolante): il gap logistico a $X_r=0.5$ è $8.0\pm0.5\%$ contro $11.8$--$12.2\%$ di misto, residenziale e suburbano ($18.4\%$ nelle villette), e a quella frequenza CW usa in media $2.03\pm0.42$ camion in più di Greedy. A $\gamma=0.9$ la soluzione è quella di $\gamma=0.5$ e, poiché il disservizio residuo è piccolo rispetto al costo (con $k_{\mathrm{scala}}\simeq1$, il termine di insoddisfazione pesa circa la metà di $F_{\mathrm{tot}}$, contro l'$80\%$ di Fabriano), la diluizione è debole e il gap resta alto: il rapporto tra i gap a $\gamma=0.9$ e $0.5$ ($0.52$) coincide con quello previsto da $(1-\gamma)\Delta F_{\mathrm{costi}}/F_{\mathrm{tot}}$.

\paragraph{Origine del vantaggio logistico per scenario.}\mbox{}\\
A pari frequenza (regime a tempo) il gap logistico è $13.4\pm1.4\%$ (grandi condomini), $16.3\pm1.3\%$ (misto), $20.9\pm1.5\%$ (residenziale), $21.9\pm1.4\%$ (suburbano) e $18.4\pm1.5\%$ (villette). La componente di chilometri è tra $11.5$ e $16.7$ punti e quella di camion tra $1.8$ e $5.6$ punti. La quota di viaggio di CW cresce da $28.6\%$ (grandi condomini) a $48.7\%$ (villette), ma non ordina da sola il gap: le villette hanno la quota maggiore e un gap logistico inferiore a residenziale e suburbano perché $\Delta V$ è piccolo ($0.83$ camion contro $1.6$ e $2.5$).

\paragraph{Sintesi.}\mbox{}\\
La struttura urbanistica determina il carico, le frequenze e il regime attivo, quindi i costi e il numero di camion; sposta invece poco il vantaggio di CW rispetto a Greedy, che a $\gamma=0.5$ resta intorno al $12\%$ (misto $9.7\%$) e dipende dal profilo in $\gamma$ solo nei grandi condomini.



\section{Benchmark 3 - Analisi della distribuzione geografica degli utenti}\label{se:b3}

\subsection{Descrizione degli scenari}

Il terzo benchmark isola l'effetto della posizione degli utenti sulla rete, mantenendo invariati numero di utenti ($N=400$), tipologia urbanistica (\textit{residenziale}), rete di base, parametri e algoritmi. Le due configurazioni sono:
\begin{itemize}
	\item \textbf{Uniforme}: i $400$ utenti sono estratti a caso tra i nodi della rete di base, quindi distribuiti su tutta l'area; coincide con lo scenario \textit{residenziale} del Benchmark 2.
	\item \textbf{Aggregata (cluster)}: i $400$ utenti sono i nodi più vicini al centro della rete, dove si trova il deposito; la distribuzione ha quindi un solo gruppo, centrato sul deposito.
\end{itemize}
Le due configurazioni differiscono quindi per la distanza tra utenti e deposito e tra utenti, non per il carico di rifiuti. Come negli altri benchmark, i risultati sono medie su 30 seed e le tre politiche ($\gamma=0.10$, $0.50$, $0.90$) sono applicate in post-elaborazione sulla stessa griglia (Sezione~\ref{se:metodo-analisi}); il confronto si basa sulla soluzione fisica e sul gap, non sul valore assoluto di $F_{\mathrm{tot}}$.

\newpage
\subsection{Discussione dei risultati}\label{sse:b3-discussione}

\begin{table}[htbp]
	\centering
	\caption{Benchmark 3: soluzione selezionata a $\gamma=0.5$ per distribuzione spaziale (media su 30 seed; somma sui cinque rifiuti, grandezze non pesate).}
	\label{tab:b3_fisica}
	\resizebox{\textwidth}{!}{%
		\begin{tabular}{llrrrrrrr}
			\hline
			Distribuzione & Algoritmo & $V$ & $V\!\cdot\!X$ & km/sett. & ore/sett. & $F_{\mathrm{ins}}$ & $F_{\mathrm{costi}}$ (\euro) & Viaggio (\%) \\
			\hline
			Uniforme & Greedy & 18.1 & 31.2 & 2751 & 215.0 & 6003 & 7169 & 50.7 \\
			Uniforme & CW     & 16.5 & 30.9 & 1986 & 198.1 & 5250 & 6568 & 39.6 \\
			Aggregata & Greedy & 15.8 & 26.4 & 1758 & 175.4 & 5959 & 5712 & 39.8 \\
			Aggregata & CW     & 14.2 & 26.4 & 1116 & 158.5 & 5466 & 5256 & 27.8 \\
			\hline
	\end{tabular}}
\end{table}

\begin{table}[htbp]
	\centering
	\caption{Benchmark 3: gap su $F_{\mathrm{tot}}$ (media $\pm$ IC95 su 30 seed) per distribuzione e politica.}
	\label{tab:b3_gap}
	\small
	\begin{tabular}{lccc}
		\hline
		Distribuzione & $\gamma=0.1$ & $\gamma=0.5$ & $\gamma=0.9$ \\
		\hline
		Uniforme  & $6.17\pm0.48$ & $11.58\pm1.04$ & $3.69\pm0.28$ \\
		Aggregata & $6.25\pm0.60$ & $8.87\pm0.65$  & $2.79\pm0.16$ \\
		\hline
	\end{tabular}
\end{table}

\paragraph{Cosa cambia con la concentrazione degli utenti.}\mbox{}\\
A $\gamma=0.5$ la concentrazione riduce i chilometri settimanali del $36\%$ per Greedy ($1758$ contro $2751$) e del $44\%$ per CW ($1116$ contro $1986$), le ore di lavoro del $18\%$ e del $20\%$, il costo logistico grezzo del $20\%$ per entrambi, i camion per giro del $13$--$14\%$ e le attivazioni settimanali del $15\%$. Il disservizio grezzo e le frequenze ottime restano praticamente uguali (per CW a $\gamma=0.5$: carta $1.4$, indifferenziata $2.5$ contro $2.6$, organico $2.6$ contro $2.8$, plastica $2.4$ contro $2.5$): la differenza di costo non nasce da scelte di servizio diverse ma dalla geometria delle rotte. Il paradosso dei camion e la resilienza si ripetono identici nelle due distribuzioni: i camion di Greedy passano da $21.7$ a $15.8$ tra $\gamma=0.1$ e $0.5$ nell'aggregata (rapporto $1.37$) e da $23.5$ a $18.1$ nell'uniforme ($1.30$), e restano costanti tra $\gamma=0.5$ e $0.9$. Anche il costo marginale è lo stesso (rapporto tra i due tratti, per CW, $7.3$ contro $7.6$).

\paragraph{Gap: la concentrazione riduce il vantaggio di CW a $\gamma\ge0.5$.}\mbox{}\\
Il gap non cambia a $\gamma=0.1$ ($6.25\pm0.60\%$ contro $6.17\pm0.48\%$) ed è più basso nella distribuzione aggregata a $\gamma=0.5$ ($8.87\pm0.65\%$ contro $11.58\pm1.04\%$) e a $\gamma=0.9$ ($2.79\pm0.16\%$ contro $3.69\pm0.28\%$), con intervalli non sovrapposti (Tabella~\ref{tab:b3_gap}). Il vantaggio di CW non cresce quindi con la concentrazione degli utenti: diminuisce, in linea con il minor peso del viaggio (il viaggio è il $27.8\%$ del tempo di giro di CW nell'aggregata contro il $39.6\%$ nell'uniforme).

\paragraph{Profilo del gap a frequenza fissata.}\mbox{}\\
A $X_r$ fissato il gap logistico è inferiore o uguale nell'aggregata a ogni frequenza: $11.13\%$ contro $11.96\%$ a $X_r=0.5$, $15.31\%$ contro $18.28\%$ a $1.0$, $20.11\%$ contro $20.92\%$ a $1.5$, $20.25\%$ contro $20.92\%$ da $X_r=2.0$ in poi (regime a tempo). Il plateau è quasi lo stesso, ma si raggiunge più tardi ($X_r=2.0$ contro $1.5$): a parità di frequenza i camion dell'aggregata, con rotte più corte, sono più pieni (saturazione fisica di Greedy $65.5\%$ contro $58.9\%$ a $X_r=1.0$, $45.8\%$ contro $39.7\%$ a $1.5$) e meno saturi in tempo ($82.3\%$ contro $88.6\%$ nel plateau), quindi la capacità resta vincolante fino a frequenze più alte. La composizione del plateau è coerente con il peso del viaggio: la componente di chilometri è $15.4$ punti nell'aggregata contro $16.7$ nell'uniforme, quella di camion $4.9$ contro $4.2$ (con $\Delta V=1.63\pm0.41$ e $1.60\pm0.49$ camion), e le due si compensano.

\begin{figure}[htbp]
	\centering
	\includegraphics[width=0.7\textwidth]{immagini/fig_gap_profilo_B3.png}
	\caption{Benchmark 3: gap logistico tra Greedy e CW a frequenza $X_r$ comune, per distribuzione uniforme e aggregata ($N=400$, media su 30 seed, fasce IC95).}
	\label{fig:b3_profilo}
\end{figure}

La curva dell'aggregata sta sotto quella dell'uniforme a ogni frequenza ($11.13\%$ contro $11.96\%$ a $X_r=0.5$, $15.31\%$ contro $18.28\%$ a $1.0$) e raggiunge il plateau più tardi ($X_r=2.0$ contro $1.5$): è il vantaggio logistico minore, non il peso del termine logistico, a spiegare il gap totale più basso.


\paragraph{Contabilità del gap a pari politica.}\mbox{}\\
Il gap dovuto al routing, cioè quello che CW ottiene a parità di frequenze (quelle scelte da Greedy), è il prodotto di due fattori: il vantaggio logistico $L_g$ di CW a quelle frequenze, $(F_{\mathrm{costi}}^{G}-F_{\mathrm{costi}}^{CW})/F_{\mathrm{costi}}^{CW}$, e il peso $w=(1-\gamma)F_{\mathrm{costi}}^{CW}/F_{\mathrm{tot}}^{CW}$ del termine logistico nel totale di CW. A $\gamma=0.5$ il termine di routing vale $8.42$ punti nell'aggregata ($L_g=17.50\%$, $w=0.48$) e $10.86$ nell'uniforme ($L_g=22.97\%$, $w=0.47$): i pesi sono quasi uguali e la differenza ($-22\%$) viene dal vantaggio logistico, che nell'aggregata è inferiore del $24\%$. La stessa struttura si ripete a $\gamma=0.9$ ($17.27\times0.16=2.78$ contro $23.21\times0.16=3.69$ punti). A $\gamma=0.1$ il fattore logistico è più basso nell'aggregata ($9.80\%$ contro $11.78\%$) ma il peso è più alto ($0.50$ contro $0.45$), quindi il termine di routing scende di meno ($4.84$ contro $5.30$ punti) e lo spostamento di $X^*$ lo compensa, perché è più grande nell'aggregata ($1.41$ contro $0.88$ punti): i gap totali risultano uguali ($6.25$ e $6.17\%$). Il minor vantaggio logistico nell'aggregata è coerente con il minor peso del viaggio, su cui CW agisce (Tabella~\ref{tab:b3_fisica}); $L_g$ è calcolato alle frequenze di Greedy, diverse per rifiuto, e non coincide quindi con il gap logistico a $X_r$ comune del paragrafo precedente.

\paragraph{Sintesi.}\mbox{}\\
La distribuzione spaziale cambia soprattutto i costi, i chilometri e il numero di camion (circa $-20\%$ di costo con gli utenti attorno al deposito), non le frequenze ottime né la struttura del compromesso tra le politiche; riduce inoltre il vantaggio di CW a $\gamma\ge0.5$, perché il viaggio, su cui CW agisce, pesa meno.



\section{Benchmark 4 - Analisi sulla Mappa Reale di Fabriano}\label{se:b4}

\subsection{Istanza reale e lettura dei risultati}

Il quarto benchmark applica le due euristiche alla rete stradale reale di Fabriano, ricavata da OpenStreetMap. Gli utenti sono i 3932 edifici residenziali estratti dalla mappa e proiettati sull'arco stradale più vicino; la tipologia di utenza (single, famiglia, palazzina piccola, palazzina grande) è stimata dalla superficie dell'edificio, e il deposito è collocato in Via Vittorio Bachelet 15 (sede di Anconambiente S.p.A.). Gli utenti occupano un'area di $16.5\,\mathrm{km}^2$ (inviluppo convesso), con una densità di $239$ utenti$/\mathrm{km}^2$. I parametri economici e operativi sono quelli della Sezione~\ref{se:configurazione}; le quantità di rifiuto prodotte restano quelle parametriche, per cui la componente reale dell'istanza riguarda la rete, la posizione e la tipologia degli utenti.

Per il confronto con le reti sintetiche (Sezione~\ref{sse:b4-sintetico-reale}) servono l'area e la composizione tipologica degli utenti reali; la scala è il fattore $\sqrt{A_{\text{Fabriano}}/100}$ con cui si riducono le coordinate del sintetico originale, che ha $A=100\,\mathrm{km}^2$:

\begin{gather*}
	\text{Area}_{\text{Fabriano}} = 16.47 \, \text{km}^2 \implies \text{scala} = 0.4058 \\[1ex]
	\text{mix: single } 0.106, \ \text{famiglia } 0.581, \\ \text{palazzina piccola } 0.267, \ \text{palazzina grande } 0.046
\end{gather*}


Il benchmark è una singola istanza deterministica che include tutti gli utenti reali, quindi i valori riportati non hanno intervalli di confidenza; la variabilità dovuta alla scelta degli utenti è valutata nella Sezione~\ref{sse:b4-sintetico-reale} con 30 sottocampioni casuali. Seguendo il metodo della Sezione~\ref{se:metodo-analisi}, le tre politiche non sono confrontate sul valore di $F_{\mathrm{tot}}$, che cambia con i pesi e con $k_{\mathrm{scala}}$ (qui $0.546$), ma sulla soluzione fisica che producono.

\subsection{Soluzione fisica al variare di \texorpdfstring{$\gamma$}{gamma}}\label{sse:b4-fisica}

La Tabella~\ref{tab:b4_fisica} riporta, per le tre politiche, la soluzione selezionata da ciascun algoritmo, e la Tabella~\ref{tab:b4_freq} le frequenze ottime per rifiuto.

\begin{table}[htbp]
	\centering
	\caption{Benchmark 4: soluzione selezionata per politica (somma sui cinque rifiuti, grandezze non pesate; le saturazioni sono medie sui rifiuti). $V$: camion per giro; $V\!\cdot\!X$: attivazioni settimanali; $F_{\mathrm{ins}}$: disservizio grezzo; $F_{\mathrm{costi}}$: costo logistico grezzo settimanale.}
	\label{tab:b4_fisica}
	\resizebox{\textwidth}{!}{%
		\begin{tabular}{llrrrrrrr}
			\hline
			$\gamma$ & Algoritmo & $V$ & $V\!\cdot\!X$ & km/sett. & ore/sett. & $F_{\mathrm{ins}}$ & $F_{\mathrm{costi}}$ (\euro) & Sat.\ fis./temp.\ (\%) \\
			\hline
			0.10 & Greedy & 145 & 92.5  & 1615 & 485  & 211\,992 & 16\,529 & 89.6 / 69.6 \\
			0.10 & CW     & 112 & 93.0  & 1017 & 579  & 174\,227 & 17\,724 & 89.7 / 78.6 \\
			0.50 & Greedy & 93  & 161.0 & 3570 & 1224 & 47\,228  & 34\,711 & 55.8 / 95.2 \\
			0.50 & CW     & 88  & 170.0 & 2104 & 1293 & 40\,344  & 36\,025 & 56.2 / 95.3 \\
			0.90 & Greedy & 93  & 199.0 & 4377 & 1512 & 33\,460  & 43\,301 & 50.1 / 95.2 \\
			0.90 & CW     & 88  & 188.0 & 2319 & 1429 & 33\,460  & 40\,307 & 53.0 / 95.3 \\
			\hline
	\end{tabular}}
\end{table}

\begin{table}[htbp]
	\centering
	\caption{Benchmark 4: frequenze ottime $X_r^*$ (ritiri a settimana) per rifiuto e politica.}
	\label{tab:b4_freq}
	\small
	\begin{tabular}{llccccc}
		\hline
		$\gamma$ & Algoritmo & Carta & Indifferenziata & Organico & Plastica & Vetro \\
		\hline
		0.10 & Greedy        & 1.0 & 0.5 & 0.5 & 1.0 & 0.5 \\
		0.10 & CW            & 1.0 & 0.5 & 1.0 & 1.5 & 0.5 \\
		0.50 & Greedy        & 1.5 & 2.0 & 2.5 & 2.5 & 0.5 \\
		0.50 & CW            & 1.5 & 3.0 & 2.5 & 2.5 & 0.5 \\
		0.90 & Greedy e CW   & 1.5 & 3.0 & 3.5 & 2.5 & 0.5 \\
		\hline
	\end{tabular}
\end{table}

\paragraph{Cosa cambia passando da una politica all'altra.}\mbox{}\\
Da $\gamma=0.10$ a $\gamma=0.50$ le attivazioni settimanali passano da $92.5$ a $161$ per Greedy e da $93$ a $170$ per CW, i chilometri settimanali da $1615$ a $3570$ e da $1017$ a $2104$, le ore di lavoro da $485$ a $1224$ e da $579$ a $1293$. In cambio il disservizio grezzo scende da $212$\,mila a $47$\,mila (Greedy, $-78\%$) e da $174$\,mila a $40$\,mila (CW, $-77\%$). Il passo successivo, da $\gamma=0.50$ a $\gamma=0.90$, riduce il disservizio solo di un ulteriore $29\%$ per Greedy e $17\%$ per CW (a $33$\,mila per entrambi), ma aumenta il costo logistico grezzo del $24.7\%$ e dell'$11.9\%$. Il costo marginale per unità di disservizio eliminata passa, per CW, da $0.137$\,\euro\ tra $\gamma=0.1$ e $0.5$ a $0.622$\,\euro\ tra $0.5$ e $0.9$ ($4.55$ volte; per Greedy da $0.110$ a $0.624$\,\euro, $5.65$ volte). Sulla griglia completa di $\gamma$ il salto più grande è però l'ultimo: da $0.622$ a $7.9$\,\euro\ per unità al passaggio a $\gamma=0.95$ (Sezione~\ref{sse:t1-soglie}).

\paragraph{Il paradosso dei camion a $\gamma=0.10$.}\mbox{}\\
A $\gamma=0.10$ vengono usati più camion per giro che a $\gamma=0.50$ ($145$ contro $93$ per Greedy, $112$ contro $88$ per CW), pur essendo la politica che minimizza il costo logistico. Il motivo è che $V_r$ conta i camion di un singolo giro, mentre il costo è proporzionale ai passaggi settimanali. A $\gamma=0.10$ la Grid Search seleziona frequenze basse (Tabella~\ref{tab:b4_freq}); poiché $Q_{ru}=W_{rt_u}/X_r$, il carico per ritiro cresce e la capacità diventa il vincolo attivo: la saturazione fisica media è $89.6\%$ (Greedy) e $89.7\%$ (CW), contro circa il $56\%$ di $\gamma=0.50$, e per l'organico raggiunge il $99.1\%$. In questo regime il numero di carichi settimanali è fissato dal rifiuto prodotto, $V_rX_r\simeq\sum_u W_{rt_u}/(C_r\cdot\mathrm{sat}_r)$, e non dipende da $X_r$: per l'organico $V_rX_r=25$ sia con Greedy ($X_r=0.5$, $50$ camion) sia con CW ($X_r=1.0$, $25$ camion). Dimezzare la frequenza raddoppia i camion per giro ma non le attivazioni settimanali: sotto la soglia di regime il costo fisso $c_rV_rX_r$ resta al suo minimo, e ciò che continua a scendere è il lavoro, perché ogni giro visita tutti gli utenti (il tempo di servizio di un giro è, ad esempio, $6390$\,min per la carta e $7668$\,min per l'organico, identico per i due algoritmi). Le attivazioni settimanali totali sono quindi $92.5$ (Greedy) e $93$ (CW) a $\gamma=0.10$, contro $161$ e $170$ a $\gamma=0.50$, e i chilometri settimanali sono meno della metà ($1615$ contro $3570$ per Greedy).

\paragraph{Il vincolo attivo per $\gamma\ge0.5$.}\mbox{}\\
Per $\gamma\ge0.5$ la saturazione temporale è circa il $95\%$ per entrambi gli algoritmi ed è il tempo del turno a determinare la flotta: $V_r=\lceil T_r/L\rceil$, con $T_r$ minuti totali di un giro. Per la carta, ad esempio, un giro richiede $6872$\,min con CW ($14.3$ turni da $480$\,min, quindi $15$ camion) e $7313$\,min con Greedy ($15.2$, quindi $16$ camion). Qui più frequenza significa più viaggi, ma non più camion per giro: $V$ è identico a $\gamma=0.5$ e a $\gamma=0.9$ ($93$ e $88$).

\paragraph{Resilienza e scelta di $\gamma$.}\mbox{}\\
Poiché le frequenze sono discrete, la soluzione è costante a tratti in $\gamma$. Su 19 valori di $\gamma$ tra $0.05$ e $0.95$ ciascun algoritmo produce solo 9 soluzioni distinte, e la stessa soluzione resta ottima per tutto l'intervallo $\gamma\in[0.55,0.90]$ (Figura~\ref{fig:b4_intervalli}), con $188$/$199$ attivazioni, $2319$/$4377$\,km e $1429$/$1512$ ore a settimana (CW/Greedy). Tra $\gamma=0.5$ e $\gamma=0.9$ le frequenze cambiano solo per l'organico (da $2.5$ a $3.5$) e, per Greedy, per l'indifferenziata (da $2.0$ a $3.0$). Il decisore non sceglie dunque un valore continuo di $\gamma$, ma uno dei punti della frontiera. Oltre il plateau l'unico passo successivo ($\gamma=0.95$, con organico a $4.0$ e vetro a $1.0$) aumenta chilometri e ore di circa il $10\%$ ($+9.7\%$ e $+10.7\%$ per CW) e il costo logistico dell'$11.6\%$ per ridurre il disservizio solo dell'$1.8\%$ (per Greedy $+9.6\%$, $+10.6\%$, $+11.5\%$ e $-1.8\%$). La scelta pratica si riduce quindi a decidere se fermarsi alla soluzione di $\gamma=0.5$ o passare a quella del plateau, che per CW costa $+11.9\%$ di costo logistico per $-17\%$ di disservizio (Figura~\ref{fig:b4_fisica}).

\begin{figure}[H]
	\centering
	\includegraphics[width=0.70\textwidth]{immagini/fig_intervalli_B4.png}
	\caption{Benchmark 4: soluzioni distinte (punti) e intervalli di $\gamma$ per cui sono ottime. Asse $x$: costo logistico grezzo settimanale; asse $y$: disservizio grezzo.}
	\label{fig:b4_intervalli}
\end{figure}

\begin{figure}[htbp]
	\centering
	\includegraphics[width=\textwidth]{immagini/fig_fisica_B4.png}
	\caption{Benchmark 4: grandezze fisiche e frequenze ottime al variare di $\gamma$ (Greedy in blu e CW in arancione; nel pannello delle frequenze CW a linea continua e Greedy tratteggiato dove le due curve differiscono).}
	\label{fig:b4_fisica}
\end{figure}

\subsection{Confronto tra reti sintetiche e rete reale}\label{sse:b4-sintetico-reale}

Il vantaggio di CW su Fabriano è sensibilmente inferiore a quello dei benchmark sintetici: a $\gamma=0.5$ il gap su $F_{\mathrm{tot}}$ è $4.2\%$, contro $9.42\pm0.47\%$ a $N=1000$, $9.20\pm0.41\%$ nel grafo standard ($N=200$) e $21.7\%$ a $N=100$. Per individuarne l'origine si sono confrontate tre possibili cause con test mirati: la topologia della rete, il numero di utenti e l'area servita (densità).


\begin{figure}[htbp]
	\centering
	\includegraphics[width=0.7\textwidth]{immagini/fig_gap_profilo_B4.png}
	\caption{Gap logistico tra Greedy e CW a frequenza $X_r$ comune con $N=1000$: sintetico originale ($A=100$\,km$^2$), sintetico con area e mix di Fabriano (B1S) e sottocampioni di utenti reali di Fabriano (B4sub) (media su 30 seed, fasce IC95).}
	\label{fig:b4_profilo}
\end{figure}

A parità di $N$, ridurre l'area del sintetico a quella di Fabriano abbassa il plateau del gap logistico fino a portarlo vicino alla curva dei sottocampioni reali: la differenza con le reti sintetiche originali è dovuta ad area, densità e mix tipologico (che nel test cambiano insieme, paragrafo «Area e densità»), più che alla struttura della rete.

\paragraph{Topologia e chilometri.}\mbox{}\\
I chilometri di un giro scalano come $\sqrt{N A}$, con $A$ area servita. Si definisce $r=\text{km per giro}/\sqrt{N\,A}$, misurato nel regime a tempo ($X_r=4$, media sui rifiuti), in cui i chilometri non dipendono da $X_r$. La Tabella~\ref{tab:b4_r} confronta Fabriano (sottocampioni e istanza completa) con tre serie sintetiche su 30 seed: il sintetico originale su rete di base da 1000 nodi ($N\le1000$) e le due serie su rete da 4000 nodi, con area di $100\,\mathrm{km}^2$ oppure con area e mix di Fabriano.

\begin{table}[htbp]
	\centering
	\caption{Chilometri per giro normalizzati, $r=\text{km}/\sqrt{N A}$ (regime a tempo, media sui rifiuti e sui 30 seed), per CW, e rapporto tra $r$ di Greedy e di CW. Sintetici: rete di base da 1000 nodi ($A=100\,\mathrm{km}^2$) e da 4000 nodi ($A=100\,\mathrm{km}^2$, oppure area e mix di Fabriano). Fabriano: sottocampioni di utenti reali per $N<3932$, istanza completa per $N=3932$.}
	\label{tab:b4_r}
	\resizebox{\textwidth}{!}{%
		\begin{tabular}{rccccccc}
			\hline
			& \multicolumn{4}{c}{$r_{CW}$} & \multicolumn{3}{c}{Greedy/CW} \\
			$N$ & 1000 nodi & 4000 nodi & 4000 nodi, area Fabriano & Fabriano & 1000 nodi & 4000 nodi, area Fabriano & Fabriano \\
			\hline
			100  & 1.157 & 1.222 & 1.219 & 1.288 & 1.43 & 1.42 & 1.43 \\
			250  & 1.074 & 1.174 & 1.166 & 1.214 & 1.52 & 1.59 & 1.52 \\
			500  & 1.005 & 1.145 & 1.131 & 1.117 & 1.61 & 1.72 & 1.62 \\
			750  & 0.967 & 1.117 & 1.103 & 1.061 & 1.69 & 1.79 & 1.71 \\
			1000 & 0.947 & 1.099 & 1.083 & 1.016 & 1.74 & 1.84 & 1.75 \\
			2000 & --    & 1.067 & 1.048 & 0.905 & --   & 1.91 & 1.85 \\
			2500 & --    & --    & --    & 0.878 & --   & --   & 1.87 \\
			3000 & --    & 1.051 & 1.029 & 0.861 & --   & 1.95 & 1.87 \\
			3500 & --    & --    & --    & 0.848 & --   & --   & 1.89 \\
			3932 & --    & 1.048 & 1.026 & 0.838 & --   & 1.95 & 1.89 \\
			\hline
	\end{tabular}}
\end{table}

Il valore assoluto di $r$ dipende dalla rete di base. A pari $N$ e area, sulla rete da 4000 nodi $r_{CW}$ è più alto che su quella da 1000 nodi del $5.6\%$ a $N=100$ e del $16.1\%$ a $N=1000$, e decresce più lentamente: $r_{CW}\simeq1.742\,N^{-0.088}$ sulla rete da 1000 nodi, $1.495\,N^{-0.044}$ su quella da 4000 nodi e $1.525\,N^{-0.049}$ con area e mix di Fabriano (residui sotto lo $0.8\%$). La legge della rete da 1000 nodi, estrapolata a $N=3932$, darebbe $0.838$, lo stesso valore di Fabriano; la misura diretta sui sintetici a $N=3932$ dà invece $1.048$ e $1.026$, per cui quella coincidenza non va interpretata.

Il confronto più omogeneo è con il sintetico che ha la stessa area, lo stesso mix e la stessa rete di base da 4000 nodi. Rispetto a questo, su Fabriano $r_{CW}$ è più alto per $N\le250$ ($+5.7\%$ e $+4.2\%$) e più basso da $N=500$ in poi, fino al $-18.3\%$ a $N=3932$: tra $N=100$ e $N=3932$ scende del $35\%$ sulla rete reale e del $16\%$ sul sintetico. Greedy segue lo stesso andamento ($+6.3\%$ a $N=100$, $-20.8\%$ a $N=3932$). Gli utenti reali non sono distribuiti uniformemente sull'inviluppo convesso: l'indice di Clark--Evans $R$, cioè la distanza media dal vicino più prossimo divisa per quella attesa con punti uniformi alla stessa densità, vale $0.87$ sui sottocampioni con $N=100$ e scende a $0.32$ con tutti gli utenti, mentre per punti uniformi resta tra $1.00$ e $1.12$. Al crescere di $N$ emergono quindi gli addensamenti, e questo va nella direzione della decrescita più rapida di $r$. L'addensamento però non ne spiega l'ampiezza: se la lunghezza del giro scalasse con un'area efficace $A R^2$, lo scarto di $r$ rispetto al sintetico dovrebbe passare da $-22\%$ a $-68\%$ tra $N=100$ e $N=3932$, mentre quello osservato passa da $+5.7\%$ a $-18.3\%$. Non spiega nemmeno lo scarto positivo per $N\le250$, la cui causa resta non isolata. Il rapporto Greedy/CW sulla rete reale coincide con quello del sintetico da 1000 nodi per $N\le1000$ ed è inferiore del $3$--$6\%$ a quello del sintetico con area di Fabriano per $N\ge250$. In chilometri a giro, a $N=3932$ CW percorre $213$\,km contro $402$ di Greedy ($-47.0\%$), contro il $-48.6\%$ del sintetico con area di Fabriano e il $-45.7\%$ del sintetico originale: sulla rete reale il margine di routing di CW è dello stesso ordine che sul sintetico.

\paragraph{Peso economico del viaggio.}\mbox{}\\
Il tempo di servizio è identico tra algoritmi e valori di $\gamma$ (differenza massima di $0.003$\,min su 960 gruppi di soluzioni), quindi il tempo di viaggio si ricava per differenza. A $X_r=4$ (regime a tempo, media sui cinque rifiuti) il viaggio è il $6.4\%$ del tempo di giro di CW a Fabriano ($11.5\%$ per Greedy), contro il $28.0\%$ ($N=1000$) e il $59.8\%$ ($N=100$) del sintetico originale; la densità è $239$ utenti$/\mathrm{km}^2$ a Fabriano e $10$ utenti$/\mathrm{km}^2$ e $1$ utenti$/\mathrm{km}^2$ rispettivamente nelle due istanze sintetiche. Poiché i chilometri crescono come $\sqrt{N A}$ mentre servizio e camion crescono con $N$, ad alta densità la quota di costo su cui il routing può agire è piccola. Il vantaggio logistico di CW a pari $X_r$ nel regime a tempo si scompone in un termine di chilometri (che costano $0.95$\,\euro/km tra costo chilometrico e tempo) e in un termine di camion ($\Delta V\,c_r\,X_r$): è pari al $7.5\%$ a Fabriano (circa il $66\%$ dai chilometri e il $34\%$ da cinque camion in meno su $93$), al $19.1\%$ a $N=1000$ sintetico e al $42.8\%$ a $N=100$.

\paragraph{Area e densità.}\mbox{}\\
Per isolare il ruolo dell'area, il benchmark sintetico è stato ripetuto con coordinate ridotte (area di $16.5\,\mathrm{km}^2$, come Fabriano) e con la composizione tipologica reale, e confrontato con sottocampioni casuali degli utenti reali di Fabriano; ogni punto è la media di 30 seed. Per arrivare a $N=3932$ (tutti gli utenti di Fabriano) le due serie sintetiche sono generate su una rete di base da 4000 nodi, da cui sono estratti gli utenti; distanze e tempi scalano linearmente con l'area. La Tabella~\ref{tab:b4_gap} e la Figura~\ref{fig:b4_gap_N} riportano il gap a pari $N$. La rete di base incide solo a $N$ piccolo nel sintetico originale: a $N=100$ il gap è $21.7\pm2.4\%$ sulla rete da 1000 nodi del Benchmark~B1 e $15.8\pm2.1\%$ su quella da 4000 nodi (causa non isolata), mentre per $N\ge250$ le due reti differiscono di meno di $0.9$ punti; per il sintetico con area di Fabriano la differenza non supera $1.2$ punti.


\begin{table}[htbp]
	\centering
	\caption{Gap su $F_{\mathrm{tot}}$ a $\gamma=0.5$ (media $\pm$ IC95 sui 30 seed): sintetico originale ($A=100\,\mathrm{km}^2$), sintetico con area e mix tipologico di Fabriano (entrambi su rete di base da 4000 nodi), utenti reali di Fabriano campionati. A $N=3932$ (tutti gli utenti) Fabriano dà $4.21\%$ (singola istanza, senza IC). Per $N\ge2500$ i sottocampioni condividono gran parte dei 3932 utenti: gli IC95 sottostimano la variabilità.}
	\label{tab:b4_gap}
	\small
	\begin{tabular}{rccc}
		\hline
		$N$ & Sintetico originale & Sintetico con area di Fabriano & Fabriano (rete reale) \\
		\hline
		100  & $15.77\pm2.11$ & $6.24\pm0.42$ & $6.57\pm0.40$ \\
		250  & $12.98\pm1.65$ & $6.50\pm0.25$ & $6.04\pm0.31$ \\
		500  & $8.18\pm0.87$  & $6.87\pm0.53$ & $5.72\pm0.12$ \\
		750  & $9.93\pm0.45$  & $8.93\pm0.85$ & $8.76\pm0.75$ \\
		1000 & $9.09\pm0.74$  & $6.78\pm0.69$ & $8.84\pm0.22$ \\
		2000 & $9.72\pm0.30$  & $5.61\pm0.35$ & $5.56\pm0.32$ \\
		2500 & --             & --            & $4.45\pm0.12$ \\
		3000 & $9.34\pm0.27$  & $5.43\pm0.26$ & $4.05\pm0.25$ \\
		3500 & --             & --            & $4.22\pm0.12$ \\
		3932 & $8.75\pm0.19$  & $4.58\pm0.25$ & $4.21$ (istanza reale) \\
		\hline
	\end{tabular}
\end{table}

\begin{figure}[htbp]
	\centering
	\includegraphics[width=\textwidth]{immagini/fig_gap_vs_N.png}
	\caption{Gap su $F_{\mathrm{tot}}$ in funzione di $N$ per $\gamma=0.1,\,0.5,\,0.9$: sintetico originale, sintetico con area e mix di Fabriano, utenti reali campionati (barre: IC95 sui 30 seed); la stella indica Fabriano con tutti gli utenti ($N=3932$).}
	\label{fig:b4_gap_N}
\end{figure}

La riduzione del gap dovuta ad area e mix non è uniforme in $N$. A $\gamma=0.5$ il gap passa da $15.8\%$ a $6.2\%$ a $N=100$ ($-9.5$ punti), da $13.0\%$ a $6.5\%$ a $N=250$, da $9.7\%$ a $5.6\%$ a $N=2000$ e da $8.8\%$ a $4.6\%$ a $N=3932$ ($-4.2$ punti); per $N=500$--$1000$ la riduzione è invece di $1.0$--$2.3$ punti. Oltre $N=1000$ il gap del sintetico originale resta tra $8.8\%$ e $9.7\%$, mentre nel sintetico con area di Fabriano scende da $5.6\%$ ($N=2000$) a $4.6\%$ ($N=3932$) e nei sottocampioni reali da $5.6\%$ a $4.2\%$ ($N=3500$). A $N=3932$ l'istanza reale ($4.21\%$) dista $0.37$ punti dal sintetico con area di Fabriano ($4.58\pm0.25\%$, il cui IC95 non la contiene) e $4.54$ punti dal sintetico originale ($8.75\pm0.19\%$); lo stesso ordine vale a $\gamma=0.1$ ($4.30$, $3.88$, $7.75$) e a $\gamma=0.9$ ($1.54$, $1.46$, $2.85$).

Sui sette $N$ comuni ($100,250,500,750,1000,2000,3000$) il sintetico con area di Fabriano e i sottocampioni reali differiscono al più di $0.8$ punti a $\gamma=0.1$ e $0.9$ e di $2.1$ punti a $\gamma=0.5$, con il massimo a $N=1000$ ($6.8\%$ contro $8.8\%$, $2.1$ punti) e la seconda differenza a $N=3000$ ($5.4\%$ contro $4.1\%$, $1.4$ punti); il segno della differenza cambia con $N$. La rete OpenStreetMap e il grafo sintetico (triangolazione di Delaunay potata) hanno struttura molto diversa, ma a pari area e mix producono gap che differiscono al più di $2.1$ punti: la differenza tra sintetico originale e Fabriano dipende quindi da area, densità e mix tipologico più che dalla struttura della rete. Questo non rende irrilevante il grafo: nel sintetico originale il solo numero di nodi della rete di base sposta il gap di $5.9$ punti a $N=100$ ($21.7\%$ con 1000 nodi, $15.8\%$ con 4000). A $N=100$ e $N=250$ e $\gamma=0.5$ la componente dei camion è assente sia nel sintetico con area di Fabriano sia nei sottocampioni reali ($\Delta V=0$: Greedy e CW usano lo stesso numero di camion), mentre nel sintetico originale CW usa il $26.7\%$ e il $17.3\%$ di camion in meno.

La forma non monotona rispetto a $N$ accompagna la riduzione percentuale dei camion di CW. In tutte e tre le serie $N=750$ ha un gap più alto di $N=500$ ($9.9\%$ contro $8.2\%$ nel sintetico originale, $8.9\%$ contro $6.9\%$ con area di Fabriano, $8.8\%$ contro $5.7\%$ nei sottocampioni), con una riduzione dei camion di CW di $-9.4\%$, $-6.3\%$ e $-7.1\%$ contro $-4.1\%$, $-0.6\%$ e $-1.9\%$. A $N=1000$ il gap dei sottocampioni resta alto ($8.84\%$, $-11.0\%$ di camion), mentre nel sintetico con area di Fabriano scende ($6.78\%$, $-2.9\%$). Il salto tra $N=500$ e $N=750$ è comune alle tre serie, ma la sua causa non è stata isolata.

A $\gamma=0.1$ il gap dei sottocampioni non segue una tendenza monotona: $3.99$, $4.00$, $3.89$ e $3.31$ per $N=2000,2500,3000,3500$, contro $3.88$ dell'istanza completa. A $X_r=0.5$ la soluzione è nel regime a capacità, dove il numero di camion è intero, e la differenza Greedy$-$CW fluttua tra $-0.6$ e $+0.3$ camion a seconda di $N$ (circa lo $0.5\%$ della flotta); l'istanza completa rientra nell'intervallo dei sottocampioni. Il meccanismo per cui CW impacchetta peggio a $N=3500$ non è stato isolato.

Il test cambia insieme area e mix tipologico. A parità di area, nel benchmark sulle tipologie urbanistiche (Sezione~\ref{se:b2}, $N=400$) il mix sposta il gap tra $9.7\%$ e $12.4\%$ a $\gamma=0.5$, per cui il suo contributo è al più di $2.7$ punti (misurati però su un solo $N$ e sul sintetico con $A=100\,\mathrm{km}^2$). Rispetto alla riduzione totale, la quota non attribuibile al mix è di almeno $6.8$ punti a $N=100$ e $1.5$ punti a $N=3932$ ed è da attribuire a area e densità; per $N=500$--$1000$, dove la riduzione totale è di $1.0$--$2.3$ punti, il mix potrebbe da solo spiegarla e il test non separa i due contributi. Il gap piccolo su Fabriano non dipende dunque da rotte meno efficaci di CW: i suoi chilometri sono quasi dimezzati, ma pesano sul costo totale molto meno che sulle reti sintetiche originali.

\subsection{Confronto tra Greedy e Clarke-Wright e discussione}

La Tabella~\ref{tab:b4_gapdec} scompone il gap su $F_{\mathrm{tot}}$ secondo il metodo della Sezione~\ref{se:metodo-analisi}, affiancandolo alla differenza fisica tra le due soluzioni.

\begin{table}[htbp]
	\centering
	\caption{Benchmark 4: scomposizione del gap su $F_{\mathrm{tot}}$ (in \% di $F^{CW}_{\mathrm{tot}}$) e differenza fisica di CW rispetto a Greedy sulle soluzioni selezionate. Le colonne $\Delta F_{\mathrm{costi}}$ e $\Delta F_{\mathrm{ins}}$ confrontano i punti scelti dai due algoritmi, che possono avere frequenze diverse, e non il solo effetto del routing.}
	\label{tab:b4_gapdec}
	\small
	\begin{tabular}{lrrrrr}
		\hline
		$\gamma$ & Gap & da routing & da spostamento $X^*$ & $\Delta F_{\mathrm{costi}}$ (\%) & $\Delta F_{\mathrm{ins}}$ (\%) \\
		\hline
		0.10 & 3.88 & 2.94 & 0.94 & $+7.2$ & $-17.8$ \\
		0.50 & 4.21 & 4.18 & 0.03 & $+3.8$ & $-14.6$ \\
		0.90 & 1.46 & 1.46 & 0.00 & $-6.9$ & $0.0$ \\
		\hline
	\end{tabular}
\end{table}

\paragraph{Regimi di capacità e di tempo.}\mbox{}\\
A parità di frequenza CW è meno costoso di Greedy in tutte le combinazioni (rifiuto, $X_r$), con un gap logistico tra $2.6\%$ e $8.7\%$, indipendente da $\gamma$. Il gap dipende dal regime: sotto la soglia di regime, dove vincola la capacità, vale in media $4.3\%$ a $X_r=0.5$ e non c'è vantaggio nel numero di camion (a $X_r=0.5$ $\Delta V=0$ per quattro rifiuti su cinque, $+1$ per il vetro); sopra la soglia, dove vincola il tempo, i chilometri risparmiati liberano anche camion e il gap sale a $7.5\%$ (carta $8.55\%$, plastica $8.65\%$, indifferenziata $7.17\%$, organico $6.97\%$, vetro $5.93\%$, già vincolato dal tempo a $X_r=0.5$). A questi valori il costo di un giro non dipende più da $X_r$ e il costo settimanale cresce linearmente con la frequenza.

\paragraph{$\gamma=0.90$: routing puro e diluizione.}\mbox{}\\
A $\gamma=0.90$ i due algoritmi scelgono le stesse frequenze, il disservizio è identico ($33\,460$) e il gap è tutto routing: CW costa il $6.9\%$ in meno ($40\,307$ contro $43\,301$\,\euro), cioè $L_g=7.43\%$ (Greedy costa il $7.4\%$ in più). Il gap su $F_{\mathrm{tot}}$ scende all'$1.46\%$ perché l'$80\%$ di $F_{\mathrm{tot}}$ è il termine di insoddisfazione, comune ai due algoritmi: il peso del termine logistico è $w=0.197$ e $7.43\%\times0.197=1.46\%$ ($L_g$ e $w$ sono definiti nella Sezione~\ref{sse:b3-discussione}). Non si tratta quindi di una perdita di margine di routing, come si potrebbe pensare dall'alta saturazione temporale: il vantaggio logistico a $\gamma=0.9$ è paragonabile a quello degli altri regimi.

\paragraph{$\gamma=0.50$ e $\gamma=0.10$: lo spostamento di frequenza.}\mbox{}\\
A $\gamma=0.50$ il gap ($4.21\%$) è quasi tutto routing a pari frequenza ($4.18\%$, circa $2400$\,\euro\ di costo: $34\,711$ contro $32\,284$\,\euro\ alle frequenze di Greedy), perché il vantaggio logistico è $L_g=7.52\%$ e il peso del termine logistico è ancora $w=0.556$ ($7.52\%\times0.556=4.18\%$). I due algoritmi, però, scelgono soluzioni diverse: CW porta l'indifferenziata a $3.0$ ritiri contro $2.0$ e ottiene $-14.6\%$ di disservizio a fronte di $+3.8\%$ di costo. Questo spostamento pesa solo $0.03\%$ su $F_{\mathrm{tot}}$ perché, vicino all'ottimo, la funzione pesata è piatta: $F_{\mathrm{tot}}$ nasconde una differenza fisica rilevante. A $\gamma=0.10$ lo spostamento vale $0.94$ punti su $3.88$: CW sceglie organico a $1.0$ ritiri (Greedy $0.5$) e plastica a $1.5$ ($1.0$), con $-17.8\%$ di disservizio e $+7.2\%$ di costo. Anche la differenza di camion ($112$ contro $145$) è conseguenza di frequenze diverse e non di fusioni migliori: alle frequenze di Greedy CW userebbe $143$ camion, quindi $31$ dei $33$ camion di differenza derivano dallo spostamento di frequenza e solo $2$ dal routing. A pari frequenza il vantaggio logistico è $L_g=5.3\%$ e il peso $w=0.55$, da cui i $2.94$ punti di routing ($5.3\%\times0.55\simeq2.9$).

\paragraph{Perché il gap è massimo a $\gamma=0.5$.}\mbox{}\\
Il gap su $F_{\mathrm{tot}}$ è il prodotto $L_g\cdot w$ (routing a pari frequenza) più lo spostamento di $X^*$. Il vantaggio logistico $L_g$ dipende dal regime in cui cade $X^*$: $5.3\%$ a $\gamma=0.1$, dove $X^*$ sta sotto la soglia di regime e vincola la capacità, e $7.5\%$ a $\gamma=0.5$ e $7.4\%$ a $\gamma=0.9$, dove $X^*$ sta sopra e vincola il tempo. Il peso $w$ del termine logistico vale $0.55$, $0.556$ e $0.197$: a $\gamma=0.9$ il termine di insoddisfazione, comune ai due algoritmi, è l'$80\%$ di $F_{\mathrm{tot}}$ e diluisce il vantaggio. Ne risultano $2.9$, $4.18$ e $1.46$ punti di routing, a cui si aggiunge lo spostamento di $X^*$ ($0.94$, $0.03$ e $0.00$ punti), per un gap di $3.88\%$, $4.21\%$ e $1.46\%$. Il massimo cade a $\gamma=0.5$ perché lì si sommano il vantaggio logistico del regime a tempo e un peso del costo ancora alto: a $\gamma=0.1$ manca il primo fattore, a $\gamma=0.9$ il secondo. Il margine tra $\gamma=0.1$ e $\gamma=0.5$ è piccolo su Fabriano ($0.33$ punti, singola istanza), ma il massimo a $\gamma=0.5$ si ritrova nei sottocampioni per tutte le $N$ (da $5.41\%$ a $6.57\%$ a $N=100$, da $4.95\%$ a $8.84\%$ a $N=1000$). In valore assoluto il vantaggio logistico resta limitato dal peso ridotto del viaggio (Sezione~\ref{sse:b4-sintetico-reale}).

\paragraph{Conclusioni per il decisore.}\mbox{}\\
Le due scelte sono indipendenti. Sull'algoritmo: a pari frequenza CW riduce i chilometri del $47\%$ e il costo logistico di circa il $7\%$, a un tempo di calcolo maggiore ($87.6$\,s contro $16.8$\,s per l'intera griglia, con un unico insieme di processi; Sezione~\ref{sse:t2-potatura}). Sul peso: la soluzione è costante a tratti in $\gamma$ e il costo marginale della riduzione di disservizio cresce di un fattore $4.55$--$5.65$ dal tratto $0.1$--$0.5$ al tratto $0.5$--$0.9$, per poi saltare da $0.62$ a $7.9$\,\euro\ per unità all'ultimo passo ($\gamma=0.95$): la soluzione del plateau ($\gamma\in[0.55,0.90]$) è l'ultima prima del salto. I punti selezionati da CW e da Greedy a $\gamma=0.1$ e $0.5$ non sono confrontabili sul solo costo ($+7.2\%$ e $+3.8\%$ per CW), perché corrispondono a frequenze diverse: con rotte più economiche CW sposta le frequenze verso più servizio ($-17.8\%$ e $-14.6\%$ di disservizio). A pari disservizio, cioè a pari frequenze, CW costa il $6.8$--$7.0\%$ in meno e usa $5$ camion in meno su $93$ per tutti i livelli di disservizio delle soluzioni, salvo quello a $\gamma=0.05$ ($-3.9\%$, un camion in meno). Ad esempio la soluzione di CW a $\gamma=0.45$ ha lo stesso disservizio di quella di Greedy a $\gamma=0.5$ ($47\,228$) con $32\,284$\,\euro\ contro $34\,711$\,\euro\ ($-7.0\%$) e $88$ camion contro $93$. La frontiera di CW sta quindi sotto quella di Greedy (Figura~\ref{fig:b4_intervalli}), e la scelta dell'algoritmo vale circa il $7\%$ del costo logistico a servizio invariato.


\section{Benchmark 5 - Analisi sul Grafo Standard}\label{se:b5}

\subsection{Descrizione dello scenario}

Il quinto benchmark confronta Greedy e Clarke-Wright su un'istanza standardizzata, generata artificialmente e indipendente dalle caratteristiche della rete urbana di Fabriano: per ciascun seed viene costruito un grafo di $N=200$ utenti (scenario \textit{residenziale}, distribuzione uniforme) su cui entrambi gli algoritmi risolvono la stessa istanza. I risultati sono medie su 30 seed con intervallo di confidenza al $95\%$ (Sezione~\ref{se:metodo-analisi}).

Le tre politiche sono le stesse dei benchmark precedenti: \textbf{Pro-Azienda} ($\gamma=0.10$), \textbf{Neutra} ($\gamma=0.50$) e \textbf{Pro-Cittadino} ($\gamma=0.90$), applicate in post-elaborazione sulla stessa griglia di frequenze. Poiché $k_{\mathrm{scala}}$ è specifico dell'istanza, il confronto tra politiche non usa il valore assoluto di $F_{\mathrm{tot}}$ ma la soluzione fisica che ciascuna produce (camion per giro, attivazioni, chilometri, ore, disservizio non pesato, saturazioni), insieme al gap
\[
\text{Gap}\,\% = \frac{F_{\text{Greedy}}-F_{\text{CW}}}{\min(F_{\text{Greedy}},F_{\text{CW}})}\cdot100 .
\]

\subsection{Discussione dei risultati}\label{sse:b5-discussione}

\begin{table}[htbp]
	\centering
	\caption{Benchmark 5 ($N=200$, grafo standard): soluzione selezionata per politica (media su 30 seed; somma sui cinque rifiuti, grandezze non pesate).}
	\label{tab:b5_fisica}
	\small
	\begin{tabular}{llrrrrrr}
		\hline
		$\gamma$ & Algoritmo & $V$ & $V\!\cdot\!X$ & km/sett. & ore/sett. & $F_{\mathrm{ins}}$ & $F_{\mathrm{costi}}$ (\euro) \\
		\hline
		0.10 & Greedy & 13.3 & 7.2  & 597  & 41.4  & 11\,877 & 1514 \\
		0.10 & CW     & 11.7 & 7.5  & 444  & 38.2  & 10\,810 & 1422 \\
		0.50 & Greedy & 10.8 & 18.1 & 1822 & 127.1 & 2866    & 4257 \\
		0.50 & CW     & 10.0 & 18.7 & 1240 & 106.7 & 2712    & 3787 \\
		0.90 & Greedy & 10.8 & 23.9 & 2392 & 167.1 & 2222    & 5637 \\
		0.90 & CW     & 10.0 & 23.3 & 1543 & 133.3 & 2221    & 4766 \\
		\hline
	\end{tabular}
\end{table}

\paragraph{Soluzione fisica per politica.}\mbox{}\\
Dalla Tabella~\ref{tab:b5_fisica}, il passaggio da $\gamma=0.1$ a $\gamma=0.5$ moltiplica per circa $2.5$ le attivazioni settimanali ($7.2\rightarrow18.1$ per Greedy, $7.5\rightarrow18.7$ per CW) e per $3.1$ (Greedy, $597\rightarrow1822$) e $2.8$ (CW, $444\rightarrow1240$) i chilometri, e riduce il disservizio grezzo di $76\%$ (Greedy) e $75\%$ (CW). Il passo da $0.5$ a $0.9$ riduce il disservizio solo di un ulteriore $22\%$ (Greedy) e $18\%$ (CW), a fronte di $+32\%$ e $+26\%$ di costo logistico. Il costo marginale per unità di disservizio eliminata passa da $0.292$ a $1.994$\,\euro\ per CW ($6.8$ volte) e da $0.304$ a $2.140$\,\euro\ per Greedy ($7.0$ volte), come negli altri benchmark; sulla griglia di 19 valori di $\gamma$ il salto più grande del costo marginale cade a $\gamma=0.80$ (Sezione~\ref{sse:t1-soglie}).

\paragraph{Paradosso dei camion e resilienza.}\mbox{}\\
A $\gamma=0.1$ i camion per giro sono $13.3$ (Greedy) e $11.7$ (CW), contro $10.8$ e $10.0$ a $\gamma=0.5$ e $0.9$, mentre le attivazioni settimanali crescono ($7.2\rightarrow18.1\rightarrow23.9$ per Greedy): a frequenze basse il carico per ritiro rende vincolante la capacità (saturazione fisica $73.8\%$ e $70.1\%$ contro $31.6\%$ e $32.5\%$ a $\gamma=0.5$), a frequenze alte il tempo del turno (saturazione temporale $85.6\%$ per Greedy e $72.3\%$ per CW). Tra $\gamma=0.5$ e $0.9$ il numero di camion è identico, mentre le ore di lavoro passano da $127$ a $167$ per Greedy e da $107$ a $133$ per CW. A parità di politica CW usa $0.8$ camion in meno di Greedy ($-7\%$) nel regime a tempo, ma la differenza non si ripete a frequenze basse ($\Delta V$ a $X_r=0.5$ pari a $0.13\pm0.41$: nessun vantaggio sotto capacità).

\begin{figure}[htbp]
	\centering
	\includegraphics[width=0.85\textwidth]{immagini/fig_fisica_B5.png}
	\caption{Benchmark 5 ($N=200$): soluzione selezionata al variare di $\gamma$ (media su 30 seed, fasce IC95); stessi pannelli della Figura~\ref{fig:b1_fisica}.}
	\label{fig:b5_camion}
\end{figure}

Da $\gamma=0.1$ a $0.5$ i camion per giro passano da $13.3$ a $10.8$ (Greedy) e da $11.7$ a $10.0$ (CW) mentre le attivazioni salgono da $7.2$ a $18.1$; oltre $\gamma=0.5$ il numero di camion non cambia e il disservizio scende solo da $2866$ a $2222$ per Greedy.


\paragraph{Gap e variabilità tra istanze.}\mbox{}\\
Il gap a $\gamma=0.5$ è $9.20\pm0.41\%$ e a $\gamma=0.9$ $2.89\pm0.11\%$; a $\gamma=0.1$ è $8.22\pm0.79\%$, con intervallo sovrapposto a quello di $\gamma=0.5$, quindi in questo benchmark il picco a $\gamma=0.5$ non è distinguibile da $\gamma=0.1$ mentre il minimo a $\gamma=0.9$ lo è. La deviazione standard del gap a $\gamma=0.1$ è circa $2.1$ punti: con il solo seed 42 si sarebbe osservato $3.12\%$ e la politica Pro-Azienda sarebbe apparsa quella con il gap minimo insieme a $\gamma=0.9$. La media sui tre seed originali (42, 123 e 7) è $7.0\%$, contro $8.22\%$ con trenta seed.

Sul gap logistico a pari frequenza (regime a tempo) CW è migliore del $21.2\pm1.1\%$, di cui circa il $18\%$ da camion ($\Delta V=0.77\pm0.16$) e l'$82\%$ da chilometri; a $X_r=0.5$ il vantaggio è $18.2\pm1.9\%$. La scomposizione a $\gamma=0.5$ attribuisce $8.99$ punti al routing e $0.21$ allo spostamento di $X^*$; a $\gamma=0.1$ $7.53$ e $0.69$: la differenza tra i due algoritmi, in questo grafo, è soprattutto di routing a pari frequenza.

\begin{figure}[htbp]
	\centering
	\includegraphics[width=\textwidth]{immagini/fig_gap_scomposizione.png}
	\caption{Scomposizione del gap su $F_{\mathrm{tot}}$ (punti percentuali, media su 30 seed) in routing a pari frequenza e spostamento di $X^*$, per le tre politiche: Benchmark 1 al variare di $N$, Benchmark 2 per scenario, Benchmark 5.}
	\label{fig:gap_scomposizione}
\end{figure}

In tutti i benchmark il gap è quasi interamente routing a pari frequenza; lo spostamento di $X^*$ conta solo a $\gamma=0.1$ ($0.69$ punti su $8.22$ in B5, tra il $12$ e il $30\%$ del gap in B1 per $N\ge250$) e si annulla a $\gamma=0.9$, dove i due algoritmi scelgono le stesse frequenze.


\paragraph{Trade-off tra le due componenti.}\mbox{}\\
A $\gamma=0.9$ CW costa il $15.4\%$ in meno a parità di disservizio ($4766$ contro $5637$\,\euro, stesse frequenze); a $\gamma=0.5$ costa il $10.9\%$ in meno e riduce il disservizio del $5.0\%$; a $\gamma=0.1$ costa il $5.3\%$ in meno e riduce il disservizio dell'$8.9\%$. In questo benchmark CW è migliore o uguale a Greedy in entrambi gli obiettivi grezzi a ogni politica.

\paragraph{Sintesi.}\mbox{}\\
Nel grafo standard valgono gli stessi regimi osservati altrove (capacità sotto la soglia di frequenza, tempo sopra, resilienza del numero di camion tra $\gamma=0.5$ e $0.9$, costo marginale crescente); la variabilità tra istanze è però sufficiente a rendere inaffidabile ogni conclusione basata su un solo seed, in particolare sul gap a $\gamma=0.1$.


\section{Benchmark Trasversale 1 - Generalità dell'Effetto di \texorpdfstring{$\gamma$}{gamma}}\label{se:t1}

\subsection{Obiettivo e dati}\label{sse:t1-dati}

Sulla mappa di Fabriano (Sezione~\ref{sse:b4-fisica}) la soluzione è costante a tratti in $\gamma$, resta la stessa per $\gamma\in[0.55,0.90]$, ha un costo marginale crescente con un salto finale e, a $\gamma=0.10$, CW sceglie frequenze più alte di Greedy. Questo benchmark verifica se tali proprietà valgono per tutte le istanze sperimentate e ne ricava il meccanismo dalla struttura del modello.

Si usano le $23$ istanze indipendenti dei benchmark precedenti: il sintetico originale e il sintetico con area e mix di Fabriano ($N$ da $100$ a $1000$, rete di base da $1000$ nodi), i cinque scenari urbanistici, lo scenario a cluster (quello uniforme coincide con lo scenario residenziale e non viene contato due volte), l'istanza completa di Fabriano, i sottocampioni di Fabriano ($N$ da $100$ a $2000$) e il grafo standard. Ogni istanza, tranne Fabriano completa, è replicata su $30$ seed. Le politiche sono i $19$ valori $\gamma=0.05,0.10,\dots,0.95$, applicati in post-elaborazione (Sezione~\ref{se:metodo-analisi}). Poiché $N$, area e $k_{\mathrm{scala}}$ cambiano da un'istanza all'altra, i valori assoluti non vengono mediati tra istanze: per ciascuna si calcola la media sui seed e si riportano mediana, minimo e massimo sulle istanze di grandezze adimensionali (frequenze, saturazioni, quota del viaggio, variazioni percentuali rispetto a $\gamma=0.10$). Le analisi sui passi di frequenza (Sezioni~\ref{sse:t1-soglie} e~\ref{sse:t1-algoritmi}) usano le griglie complete $X_r=0.5,\dots,6.0$ delle $22$ istanze replicate.

\newpage
\subsection{Effetto di \texorpdfstring{$\gamma$}{gamma} sulla soluzione fisica}\label{sse:t1-fisica}

\begin{table}[htbp]
	\centering
	\caption{Soluzione fisica per politica sulle 23 istanze: mediana [minimo--massimo] delle medie sui seed. Frequenze in ritiri a settimana; saturazioni e quota del viaggio (sul tempo di lavoro) medie sui rifiuti; $\Delta$: variazione percentuale rispetto a $\gamma=0.10$ dello stesso algoritmo; $V$: camion per giro.}
	\label{tab:t1_fisica}
	\resizebox{\textwidth}{!}{%
		\begin{tabular}{lcccccc}
			\hline
			& \multicolumn{3}{c}{Clarke-Wright} & \multicolumn{3}{c}{Greedy} \\
			Grandezza & $\gamma=0.10$ & $\gamma=0.50$ & $\gamma=0.90$ & $\gamma=0.10$ & $\gamma=0.50$ & $\gamma=0.90$ \\
			\hline
			$X^*$ organico        & 0.87 [0.50--1.43] & 2.73 [2.50--5.00] & 3.68 [2.50--5.00] & 0.62 [0.50--0.98] & 2.50 [2.50--5.00] & 3.68 [2.50--5.00] \\
			$X^*$ indifferenziata & 0.55 [0.50--0.73] & 2.60 [2.00--4.00] & 3.00 [2.00--4.00] & 0.50 [0.50--0.53] & 2.17 [2.00--4.00] & 3.00 [2.00--4.00] \\
			$X^*$ plastica        & 1.03 [0.72--1.50] & 2.47 [1.78--3.50] & 2.50 [2.00--3.50] & 0.83 [0.50--1.00] & 2.33 [1.58--3.50] & 2.50 [2.00--3.50] \\
			$X^*$ carta           & 0.67 [0.50--1.00] & 1.43 [1.00--2.00] & 1.50 [1.00--2.00] & 0.50 [0.50--1.00] & 1.28 [1.00--2.00] & 1.50 [1.00--2.00] \\
			$X^*$ vetro           & 0.50 [0.50--0.50] & 0.50 [0.50--1.00] & 0.68 [0.50--1.00] & 0.50 [0.50--0.50] & 0.50 [0.50--1.00] & 0.68 [0.50--1.00] \\
			Sat.\ fisica (\%)     & 78.8 [39.2--94.1] & 40.3 [15.1--60.0] & 34.9 [14.8--60.0] & 82.0 [38.7--96.3] & 41.0 [14.1--58.8] & 31.4 [12.5--58.8] \\
			Sat.\ temporale (\%)  & 66.1 [34.5--81.3] & 79.6 [63.3--95.3] & 79.6 [63.3--95.3] & 70.8 [38.8--85.5] & 88.2 [60.9--95.7] & 88.2 [60.9--95.7] \\
			Quota viaggio (\%)    & 31.9 [7.0--59.8]  & 28.6 [6.4--59.8]  & 28.6 [6.4--59.8]  & 44.9 [13.3--68.3] & 40.2 [11.5--67.9] & 40.2 [11.5--67.9] \\
			$\Delta F_{\mathrm{ins}}$ (\%)   & -- & $-76.8$ [$-95.4$; $-70.4$] & $-80.7$ [$-95.4$; $-76.5$] & -- & $-78.8$ [$-95.7$; $-72.7$] & $-84.3$ [$-95.7$; $-80.3$] \\
			$\Delta F_{\mathrm{costi}}$ (\%) & -- & $+157$ [$+97$; $+192$] & $+193$ [$+127$; $+235$] & -- & $+165$ [$+110$; $+240$] & $+235$ [$+162$; $+366$] \\
			$\Delta V$ (\%)                  & -- & $-17.4$ [$-67.4$; $0.0$] & $-17.4$ [$-67.4$; $0.0$] & -- & $-25.3$ [$-69.7$; $0.0$] & $-25.3$ [$-69.7$; $0.0$] \\
			$\Delta$ km (\%)                 & -- & $+172$ [$+104$; $+288$] & $+216$ [$+128$; $+288$] & -- & $+179$ [$+121$; $+294$] & $+255$ [$+171$; $+358$] \\
			$\Delta$ ore (\%)                & -- & $+175$ [$+113$; $+363$] & $+224$ [$+140$; $+363$] & -- & $+186$ [$+145$; $+402$] & $+271$ [$+200$; $+402$] \\
			\hline
	\end{tabular}}
\end{table}

La Tabella~\ref{tab:t1_fisica} conferma su tutte le istanze l'andamento osservato a Fabriano. Da $\gamma=0.10$ a $\gamma=0.50$ il disservizio scende in mediana del $76.8\%$ per CW (tra il $70.4\%$ e il $95.4\%$ a seconda dell'istanza) e del $78.8\%$ per Greedy, mentre il costo logistico cresce del $157\%$ e del $165\%$, con chilometri e ore di lavoro che crescono nella stessa misura. Il passo da $\gamma=0.50$ a $\gamma=0.90$ aggiunge una riduzione di soli $4$--$6$ punti del disservizio (riferita al livello di $\gamma=0.10$) contro $36$--$70$ punti di costo. I camion per giro, invece, diminuiscono ($-17\%$ per CW, $-25\%$ per Greedy) e sono gli stessi a $\gamma=0.50$ e $\gamma=0.90$ su tutta la distribuzione, così come la saturazione temporale e la quota del viaggio.

L'invarianza dipende dal regime del vincolo. Sotto una frequenza di soglia $\bar X$ (ricavata nella Sezione~\ref{sse:t2-regimi}) il carico per ritiro $Q_{ru}=W_{rt_u}/X_r$ è grande e la capacità determina il numero di camion; sopra $\bar X$ il vincolo attivo è il turno e il giro non cambia più con $X_r$. Nei sottocampioni di Fabriano con $N=1000$ (organico, CW, 30 seed), da $X_r=0.5$ a $X_r=2$ i camion per giro passano da $13.3$ a $5.0$; da $X_r=2$ in poi ogni giro usa $5.0$ camion, percorre $130.2$\,km, dura $2264$\,min e satura il turno al $94.3\%$, mentre la saturazione fisica scende come $1/X_r$ (dal $62.9\%$ a $X_r=2$ al $21.0\%$ a $X_r=6$). Il costo settimanale è quindi lineare in $X_r$: ogni mezzo ritiro settimanale in più costa la stessa cifra ($606$\,\euro{} in questo esempio). Tra $\gamma=0.50$ e $\gamma=0.90$ cambiano solo frequenze già sopra la soglia (organico e indifferenziata; il vetro resta in mediana a $0.5$), quindi cambia il numero di giri e non il giro. La saturazione fisica del $31$--$35\%$ a $\gamma=0.90$ non indica rotte inefficienti: il camion termina il turno prima di riempirsi perché ritira meno rifiuto a ogni passaggio.

Tra i due algoritmi la differenza fisica sta nel giro: nel regime a tempo Greedy satura il turno all'$88.2\%$ contro il $79.6\%$ di CW (mediane), con una quota di viaggio del $40.2\%$ contro il $28.6\%$.

\subsection{Struttura del disservizio: la frequenza di minimo disservizio}\label{sse:t1-quantile}

L'insoddisfazione di un rifiuto, $F_{\mathrm{insoddis}}(X_r)=\sum_u\bigl[\alpha\max(0,x^*_{ru}-X_r)+\beta\max(0,X_r-x^*_{ru})\bigr]$, è convessa e lineare a tratti in $X_r$, con pendenza
\begin{equation}\label{eq:t1_pendenza}
	\frac{dF_{\mathrm{insoddis}}}{dX_r}=-\alpha\,n_>(X_r)+\beta\,n_<(X_r),
\end{equation}
dove $n_>(X_r)$ e $n_<(X_r)$ sono gli utenti con frequenza ideale rispettivamente maggiore e minore di $X_r$. La pendenza cambia segno quando la quota di utenti ancora sotto-serviti scende sotto $\beta/(\alpha+\beta)=1/6$: la frequenza di minimo disservizio è il quantile $\alpha/(\alpha+\beta)=83.3\%$ delle frequenze ideali degli utenti. Poiché $x^*_{ru}$ è il prodotto di un valore base per rifiuto e di un moltiplicatore per tipologia ($0.7$, $1.0$, $1.5$, $2.0$; Tabella~\ref{tab:xstar_servizio}), il quantile coincide con la frequenza ideale di una sola tipologia e dipende soltanto dal mix degli utenti. La Tabella~\ref{tab:t1_quantile} riporta il quantile, il punto della griglia con disservizio minimo e, per confronto, la frequenza scelta a $\gamma=0.95$: coincidono in tutti i $30$ casi.

\begin{table}[htbp]
	\centering
	\caption{Frequenza ideale al quantile $83.3\%$ $\rightarrow$ frequenza della griglia con disservizio minimo, calcolata con le proporzioni nominali del mix. In tutti i 30 casi coincide con $X_r^*$ di CW a $\gamma=0.95$ (mediana sui seed; Fabriano: istanza completa).}
	\label{tab:t1_quantile}
	\resizebox{\textwidth}{!}{%
		\begin{tabular}{llccccc}
			\hline
			Scenario & Tipologia al quantile & Organico & Indifferenziata & Plastica & Carta & Vetro \\
			\hline
			Residenziale     & palazzina piccola & 3.75$\to$4.0 & 3.00$\to$3.0 & 2.55$\to$2.5 & 1.50$\to$1.5 & 0.75$\to$1.0 \\
			Suburbano        & palazzina piccola & 3.75$\to$4.0 & 3.00$\to$3.0 & 2.55$\to$3.0 & 1.50$\to$1.5 & 0.75$\to$1.0 \\
			Misto periferia  & palazzina piccola & 3.75$\to$4.0 & 3.00$\to$3.0 & 2.55$\to$3.0 & 1.50$\to$1.5 & 0.75$\to$1.0 \\
			Villette         & famiglia          & 2.50$\to$2.5 & 2.00$\to$2.0 & 1.70$\to$2.0 & 1.00$\to$1.0 & 0.50$\to$0.5 \\
			Grandi condomini & palazzina grande  & 5.00$\to$5.0 & 4.00$\to$4.0 & 3.40$\to$3.5 & 2.00$\to$2.0 & 1.00$\to$1.0 \\
			Fabriano         & palazzina piccola & 3.75$\to$4.0 & 3.00$\to$3.0 & 2.55$\to$2.5 & 1.50$\to$1.5 & 0.75$\to$1.0 \\
			\hline
	\end{tabular}}
\end{table}

Dalla~\eqref{eq:t1_pendenza} si ottiene anche quanto vale ciascun mezzo ritiro in più. Il primo passo dell'organico, da $0.5$ a $1.0$, trova tutti gli utenti sotto-serviti e riduce il disservizio di $\alpha/2=5$ unità per utente. A Fabriano il quantile dell'organico ($3.75$) e quello del vetro ($0.75$) cadono a metà di una cella della griglia: nel passo da $3.5$ a $4.0$ (e da $0.5$ a $1.0$ per il vetro) nel primo quarto della cella sono ancora sotto-serviti il $31.3\%$ degli utenti (palazzine piccole e grandi) e il disservizio scende, nel secondo resta sotto-servito solo il $4.6\%$ (palazzine grandi) e il disservizio risale; il guadagno netto è di $0.077$ unità per utente, $65$ volte meno del primo passo. Nel mix residenziale lo stesso passo vale $0.2$ unità per utente ($189$ unità misurate sul sintetico con $N=1000$). Nelle villette e nei grandi condomini, invece, il quantile cade su un punto della griglia e questo passo a guadagno quasi nullo non esiste.

\subsection{Soglie di \texorpdfstring{$\gamma$}{gamma}, soluzioni discrete e plateau}\label{sse:t1-soglie}

Il problema è separabile per rifiuto e $X_r^*(\gamma)$ si muove lungo l'inviluppo convesso inferiore dei punti $(F_{\mathrm{costi}}(X_r),F_{\mathrm{insoddis}}(X_r))$. Il passo da $X_a$ a $X_b$ conviene quando $\gamma\,k_{\mathrm{scala}}\,(F_a-F_b)\ge(1-\gamma)(C_b-C_a)$, cioè per
\begin{equation}\label{eq:t1_soglia}
	\gamma\ge\gamma^*=\frac{s}{k_{\mathrm{scala}}+s},\qquad s=\frac{C_b-C_a}{F_a-F_b},
\end{equation}
dove $s$ è il costo in euro di un'unità di disservizio evitata. Ogni passo ha la sua soglia, il numero di passi è finito e la soluzione è costante tra due soglie consecutive: sui 19 valori di $\gamma$ le soluzioni distinte sono in mediana $9$ per CW e $8$ per Greedy (da $6$ a $10$ a seconda dell'istanza), e l'intervallo più lungo con la stessa soluzione va per CW in mediana da $\gamma=0.55$ ($0.35$--$0.60$) a $\gamma=0.90$ ($0.75$--$0.95$), $7$ passi della griglia su $19$ ($5$--$13$).

Le soglie si separano in due gruppi. I passi di avvicinamento al quantile hanno guadagni grandi e quindi $s$ piccolo; i passi che attraversano il quantile hanno un guadagno quasi nullo (Sezione~\ref{sse:t1-quantile}) ma costano un mezzo giro intero, perché cadono nel regime a tempo, e quindi hanno $s$ grande e $\gamma^*$ vicino a $1$. Il plateau è il vuoto tra i due gruppi. Su Fabriano l'ultimo passo di avvicinamento della frontiera di CW costa $s=0.622$\,\euro{} per unità e, con $k_{\mathrm{scala}}=0.546$, entra a $\gamma^*=0.533$; il passo successivo costa $s=7.90$\,\euro{} ed entra a $\gamma^*=0.935$. La~\eqref{eq:t1_soglia} riproduce quindi esattamente il plateau $[0.55,0.90]$ della Figura~\ref{fig:b4_intervalli}. Sulle 22 istanze replicate, in almeno $406$ coppie istanza-seed su $660$ ($62\%$, CW) i passi che stanno sopra il vuoto più ampio tra soglie sono proprio quelli che attraversano il quantile, organico $3.5\to4.0$ e vetro $0.5\to1.0$.

\begin{figure}[htbp]
	\centering
	\includegraphics[width=\textwidth]{immagini/fig_t1_frontiere.png}
	\caption{Frontiere normalizzate su sei istanze: costo logistico e disservizio di ogni soluzione divisi per quelli di Greedy a $\gamma=0.05$ (media sui seed; un punto per valore di $\gamma$). Greedy tratteggiato, CW continuo; le etichette indicano la soluzione di CW a $\gamma=0.10$, $0.50$ e $0.90$.}
	\label{fig:t1_frontiere}
\end{figure}

La forma della frontiera ne è la conseguenza (Figura~\ref{fig:t1_frontiere}). Il costo marginale, in euro per unità di disservizio evitata, cresce lungo la frontiera in tutte le coppie istanza-seed: il primo passo costa in mediana $0.087$\,\euro{} per unità ($0.050$--$0.156$), l'ultimo $4.95$\,\euro{} ($0.49$--$9.10$), con un rapporto mediano di $56$ ($4$--$158$). In $17$ istanze su $23$ il salto più grande cade tra $\gamma=0.85$ e $\gamma=0.95$; a Fabriano è il passo che attraversa il quantile. Le eccezioni più nette sono le villette ($\gamma=0.50$) e i grandi condomini ($\gamma=0.35$), dove il quantile cade sulla griglia: lì la frontiera termina senza il passo costoso e $\gamma=0.50$ e $\gamma=0.90$ danno la stessa soluzione nel $57\%$ e nel $100\%$ dei seed (CW). Nelle altre quattro istanze (sottocampioni e sintetico con mix di Fabriano a $N=100$, cluster, grafo standard) il salto cade tra $\gamma=0.65$ e $\gamma=0.80$.

Ne seguono tre indicazioni per la scelta di $\gamma$. La prima è che il decisore sceglie tra poche soluzioni discrete e non tra valori continui di $\gamma$: tutti i valori del plateau, in mediana da $0.55$ a $0.90$, producono lo stesso servizio. La seconda è che lo stesso valore di $\gamma$ non ha lo stesso significato in istanze diverse: per la~\eqref{eq:t1_soglia} un passo con costo $s$ entra a $\gamma^*=s/(k_{\mathrm{scala}}+s)$, e $k_{\mathrm{scala}}$ varia tra $0.55$ e $2.05$. La grandezza confrontabile è il prezzo che si accetta di pagare per unità di disservizio evitata: in mediana $0.087$\,\euro{} per unità per il primo passo e $4.95$\,\euro{} (tra $0.5$ e $9$) per l'ultimo. La terza è un limite del modello: la posizione del plateau e del salto finale dipende dal passo della griglia ($0.5$), che decide se il quantile cade dentro una cella, e dal rapporto $\alpha/\beta$, che fissa il quantile. Una griglia più fine dividerebbe il passo finale in due passi con costi diversi.

\subsection{Confronto tra Greedy e Clarke-Wright}\label{sse:t1-algoritmi}

\begin{table}[htbp]
	\centering
	\caption{Confronto tra algoritmi. A pari $\gamma$: quota di coppie istanza-seed in cui la soluzione di CW domina quella di Greedy (costo e disservizio non peggiori, almeno uno migliore) o le due non sono confrontabili (CW ha meno disservizio e più costo); Greedy domina solo nello $0.5\%$ dei casi a $\gamma=0.05$ (661 coppie, Fabriano completa inclusa). A pari frequenze: le frequenze scelte da Greedy vengono servite con le rotte di CW (stesso disservizio); quota di casi in cui CW costa meno e risparmio di costo mediano [minimo--massimo] (660 coppie).}
	\label{tab:t1_algoritmi}
	\small
	\begin{tabular}{rcccc}
		\hline
		& \multicolumn{2}{c}{A pari $\gamma$} & \multicolumn{2}{c}{Alle frequenze di Greedy} \\
		$\gamma$ & CW domina (\%) & Non confrontabili (\%) & CW costa meno (\%) & Risparmio (\%) \\
		\hline
		0.05 & 99.2  & 0.3  & 99.6  & 8.9 [$-1.2$; 25.9] \\
		0.10 & 52.0  & 48.0 & 99.7  & 9.3 [$-1.2$; 25.9] \\
		0.15 & 41.5  & 58.5 & 100.0 & 11.5 [2.4; 29.9] \\
		0.20 & 82.6  & 17.4 & 100.0 & 13.0 [0.2; 36.1] \\
		0.30 & 99.5  & 0.5  & 100.0 & 13.2 [0.7; 36.8] \\
		0.45 & 94.3  & 5.7  & 100.0 & 13.2 [0.7; 36.8] \\
		0.50 & 86.1  & 13.9 & 100.0 & 13.0 [1.2; 37.1] \\
		0.55 & 90.2  & 9.8  & 100.0 & 13.1 [0.3; 37.3] \\
		0.70 & 99.8  & 0.2  & 100.0 & 13.2 [0.3; 37.3] \\
		0.90 & 99.8  & 0.2  & 100.0 & 13.3 [0.3; 36.3] \\
		\hline
	\end{tabular}
\end{table}

A pari frequenza il disservizio non dipende dalle rotte, e infatti è identico per i due algoritmi in tutte le $39\,600$ combinazioni (istanza, seed, rifiuto, $X_r$). CW costa meno nel $95.3\%$ dei casi a $X_r=0.5$, nel $98.2\%$ a $X_r=1.0$ e nel $99.2$--$99.3\%$ per $X_r\ge1.5$. Il vantaggio sotto capacità non viene dai camion: a $X_r=0.5$ CW usa più camion di Greedy nel $17.8\%$ dei casi e meno nel $7.8\%$, mentre per $X_r\ge1.5$ ne usa meno nel $28$--$30\%$ e più solo nell'$1.5$--$2.2\%$. Servendo con le rotte di CW le stesse frequenze scelte da Greedy, cioè a parità di disservizio, CW costa meno in tutti i casi per $\gamma\ge0.15$ e nel $99.6$--$99.7\%$ per $\gamma\le0.10$, con un risparmio mediano tra l'$8.9\%$ e il $13.3\%$ (Tabella~\ref{tab:t1_algoritmi}). Le cinque eccezioni sono tutte nel sintetico con mix di Fabriano ($N=250$ e $500$, $\gamma\le0.10$), cioè nel regime a capacità, con un extra-costo massimo dell'$1.2\%$. La frontiera di CW sta quindi sotto quella di Greedy (Figura~\ref{fig:t1_frontiere}).

A pari $\gamma$ il confronto è diverso, perché i due algoritmi possono scegliere frequenze diverse. Fuori da due bande CW domina in almeno il $97\%$ dei casi (oltre il $99\%$ per $\gamma\ge0.65$). Nelle bande $\gamma=0.10$--$0.20$ e $\gamma=0.45$--$0.55$, invece, in una quota fino al $58.5\%$ dei casi le due soluzioni non sono confrontabili: CW ha meno disservizio ma costa di più. Il meccanismo è la~\eqref{eq:t1_soglia}. Per uno stesso passo di frequenza il guadagno di disservizio è identico, mentre il costo del passo è più basso per CW (in mediana del $9.6$--$14.4\%$ a seconda della frequenza di partenza), quindi la soglia $\gamma^*$ di CW è più bassa: lo è nel $98.6\%$ dei $9979$ passi presenti in entrambi gli algoritmi, in mediana di $0.016$--$0.029$. Per $\gamma$ compreso tra le due soglie CW ha già aumentato la frequenza e Greedy no. Le due bande sono quelle in cui cadono le soglie dei passi di avvicinamento; nel plateau non ci sono soglie e i due algoritmi scelgono le stesse frequenze. Questa è la forma generale del paradosso dei camion osservato a Fabriano (Sezione~\ref{sse:b4-fisica}): a $\gamma=0.10$ le frequenze di CW sono più alte su tutte le istanze (organico $0.87$ contro $0.62$, plastica $1.03$ contro $0.83$, carta $0.67$ contro $0.50$ in mediana, Tabella~\ref{tab:t1_fisica}), e la differenza di camion e costi tra le due soluzioni dipende da queste frequenze, non dalle rotte.

\subsection{Sintesi}\label{sse:t1-sintesi}

Le proprietà osservate a Fabriano valgono per tutte le 23 istanze e discendono dalla struttura del modello. Oltre una frequenza di soglia il giro è fisso e il costo cresce linearmente con il numero di giri, mentre il guadagno di disservizio di ogni mezzo ritiro diminuisce fino ad annullarsi al quantile $83.3\%$ delle frequenze ideali: la frontiera è convessa e ogni passo di frequenza entra a una soglia $\gamma^*=s/(k_{\mathrm{scala}}+s)$. Il plateau di $\gamma$ è il vuoto tra le soglie dei passi di avvicinamento e quelle dei passi che attraversano il quantile, e il salto finale del costo marginale nasce dal fatto che il quantile cade a metà di una cella della griglia. CW domina Greedy a parità di disservizio, con un risparmio di costo mediano del $9$--$13\%$; a parità di $\gamma$ la dominanza vale fuori da due bande, dentro le quali CW anticipa ogni aumento di frequenza perché ogni passo gli costa meno.

\section{Benchmark Trasversale 2 - Profilo in Frequenza e Ricerca Potata}\label{se:t2}

\subsection{Obiettivo e dati}\label{sse:t2-dati}

Il Benchmark Trasversale 1 ha studiato la frequenza scelta al variare di $\gamma$. Questo benchmark guarda invece l'intera griglia $X_r=0.5,1.0,\dots,6.0$, cioè come cambiano camion, saturazioni e costi a frequenza fissata, che sono grandezze indipendenti da $\gamma$ (Sezione~\ref{se:metodo-analisi}). Si usano le stesse 23 istanze di T1, per un totale di $6610$ profili (istanza, seed, algoritmo, rifiuto) di 12 punti ciascuno. I valori di $F_{\mathrm{tot}}$ si confrontano solo all'interno dello stesso profilo, mai tra istanze o tra valori di $\gamma$. L'analisi porta a una proposta di miglioramento della procedura di ricerca, verificata su 29 istanze ricalcolate (Sezione~\ref{sse:t2-potatura}).

\subsection{Regime a capacità e regime a tempo}\label{sse:t2-regimi}

Sia $W_r$ il rifiuto prodotto in una settimana dagli utenti e $C_r$ la capacità di un camion. Per un ritiro servono almeno $W_r/(X_r C_r)$ camion per la capacità e $V^{\mathrm{t}}_r$ camion per il turno, dove $V^{\mathrm{t}}_r$ è il numero di camion del regime a tempo, che non dipende da $X_r$ (Sezione~\ref{sse:t1-fisica}). Il regime a tempo inizia quando la capacità smette di essere vincolante, cioè per
\begin{equation}\label{eq:t2_soglia}
	X_r\ge\bar X_r=\frac{W_r}{C_r\,V^{\mathrm{t}}_r}=X_r\,\frac{\text{sat. fisica}(X_r)}{100}\quad(X_r\text{ nel regime a tempo}),
\end{equation}
dove l'ultima uguaglianza segue da $\text{sat. fisica}=W_r/(X_r V^{\mathrm{t}}_r C_r)$ e permette di calcolare la soglia da qualunque punto del regime a tempo (qui $X_r=6$). La Tabella~\ref{tab:t2_soglia} confronta la soglia prevista con quella osservata, cioè la frequenza da cui il numero di camion per giro resta costante. Arrotondata al punto successivo della griglia, la previsione coincide con l'osservazione nel $95.8\%$ dei profili e non la anticipa mai; nel $4.2\%$ restante la soglia osservata arriva più tardi perché il carico non entra perfettamente nei camion previsti (con CW più spesso che con Greedy: fino al $10.7\%$ dei profili dell'organico contro il $5.0\%$).

\newpage
\begin{table}[htbp]
	\centering
	\caption{Soglia tra regime a capacità e regime a tempo per rifiuto (6610 profili): soglia osservata e prevista dalla~\eqref{eq:t2_soglia}, mediana [minimo; massimo]; quota di profili in cui la soglia osservata coincide con la previsione arrotondata alla griglia o arriva dopo; saturazione fisica alla soglia osservata.}
	\label{tab:t2_soglia}
	\resizebox{\textwidth}{!}{%
		\begin{tabular}{llcccccc}
			\hline
			Rifiuto & Algoritmo & Soglia osservata & Soglia prevista & Coincide (\%) & Dopo (\%) & Sat.\ fisica alla soglia (\%) \\
			\hline
			Organico        & CW     & 1.5 [0.5; 3.5] & 1.00 [0.25; 2.84] & 89.3  & 10.7 & 74 [48; 100] \\
			& Greedy & 1.0 [0.5; 2.5] & 0.95 [0.25; 2.34] & 95.0  & 5.0  & 72 [50; 100] \\
			Indifferenziata & CW     & 1.0 [0.5; 2.0] & 0.62 [0.16; 1.77] & 95.9  & 4.1  & 69 [31; 100] \\
			& Greedy & 1.0 [0.5; 1.5] & 0.60 [0.16; 1.46] & 99.5  & 0.5  & 66 [31; 100] \\
			Plastica        & CW     & 1.0 [0.5; 3.0] & 0.91 [0.22; 2.49] & 89.6  & 10.4 & 73 [44; 99]  \\
			& Greedy & 1.0 [0.5; 3.0] & 0.86 [0.21; 2.49] & 96.7  & 3.3  & 73 [42; 100] \\
			Carta           & CW     & 1.0 [0.5; 2.5] & 0.71 [0.17; 1.95] & 93.9  & 6.1  & 72 [35; 100] \\
			& Greedy & 1.0 [0.5; 2.5] & 0.68 [0.17; 1.95] & 98.0  & 2.0  & 75 [33; 100] \\
			Vetro           & CW     & 0.5 [0.5; 1.0] & 0.25 [0.06; 0.70] & 100.0 & 0.0  & 51 [12; 75]  \\
			& Greedy & 0.5 [0.5; 1.0] & 0.24 [0.06; 0.69] & 100.0 & 0.0  & 48 [13; 69]  \\
			\hline
	\end{tabular}}
\end{table}

\begin{figure}[htbp]
	\centering
	\includegraphics[width=\textwidth]{immagini/fig_t2_regimi.png}
	\caption{Fabriano, istanza completa, CW: saturazione temporale e fisica in funzione di $X_r$ per rifiuto. La linea punteggiata è la soglia prevista dalla~\eqref{eq:t2_soglia}; nel titolo, la frequenza da cui il numero di camion per giro resta costante.}
	\label{fig:t2_regimi}
\end{figure}

La Figura~\ref{fig:t2_regimi} mostra il meccanismo su Fabriano. Sotto la soglia la saturazione fisica è vicina al $100\%$ e quella temporale cresce con $X_r$, perché a ogni ritiro c'è meno rifiuto e lo stesso camion visita più utenti; sopra la soglia la saturazione temporale resta ferma intorno al $95\%$ e quella fisica scende come $1/X_r$. Per indifferenziata, carta e vetro la soglia osservata è il primo punto della griglia dopo quella prevista, per organico e plastica il secondo. Il vetro, che ha la capacità più alta (a pari merito con l'indifferenziata) e la produzione per utente più bassa, è nel regime a tempo già a $X_r=0.5$ in quasi tutti i profili (soglia prevista $0.25$ in mediana): ogni mezzo ritiro in più del vetro costa quindi un giro intero, il che contribuisce a rendere il suo passo da $0.5$ a $1.0$ uno degli ultimi a entrare nella soluzione (Sezione~\ref{sse:t1-soglie}).

\subsection{Sensibilità attorno alla frequenza ottima}\label{sse:t2-sensibilita}

Sul singolo profilo $F_{\mathrm{tot}}(X_r)$ è unimodale sulla griglia nel $99.8\%$ dei casi (sui $6610$ profili e 19 valori di $\gamma$; $97.8\%$ a $\gamma=0.10$), ma il minimo non cade in generale nel punto di incrocio tra la componente di disservizio e quella di costo: cade dove il costo di un mezzo ritiro in più eguaglia il disservizio che evita, secondo la soglia~\eqref{eq:t1_soglia}. La Tabella~\ref{tab:t2_sensibilita} mostra che cosa succede spostandosi di mezzo ritiro dall'ottimo a $\gamma=0.5$.

\begin{table}[htbp]
	\centering
	\caption{Effetto di mezzo ritiro in più o in meno rispetto a $X_r^*$ (CW, $\gamma=0.5$): variazioni percentuali di costo logistico, disservizio e $F_{\mathrm{tot}}$, mediana [minimo; massimo] sulle 23 istanze (media sui seed). Per il vetro il passo in meno esiste solo in 2 istanze.}
	\label{tab:t2_sensibilita}
	\resizebox{\textwidth}{!}{%
		\begin{tabular}{llccc}
			\hline
			Rifiuto & Passo & $\Delta F_{\mathrm{costi}}$ (\%) & $\Delta F_{\mathrm{ins}}$ (\%) & $\Delta F_{\mathrm{tot}}$ (\%) \\
			\hline
			Organico        & $+0.5$ & $+18.7$ [$+10.0$; $+20.0$] & $-14.4$ [$-18.9$; $+69.6$] & $+2.8$ [$+1.2$; $+28.4$] \\
			& $-0.5$ & $-18.7$ [$-20.0$; $-10.0$] & $+82.7$ [$+37.1$; $+217.0$] & $+26.3$ [$+1.9$; $+35.6$] \\
			Indifferenziata & $+0.5$ & $+20.0$ [$+12.5$; $+25.0$] & $+11.9$ [$-4.9$; $+87.0$] & $+13.3$ [$+6.8$; $+36.8$] \\
			& $-0.5$ & $-20.0$ [$-25.0$; $-12.5$] & $+67.3$ [$+44.7$; $+271.2$] & $+18.0$ [$+0.1$; $+48.9$] \\
			Plastica        & $+0.5$ & $+20.3$ [$+14.3$; $+28.6$] & $+17.0$ [$-7.6$; $+84.9$] & $+18.4$ [$+9.1$; $+25.1$] \\
			& $-0.5$ & $-20.3$ [$-28.6$; $-14.3$] & $+49.9$ [$+34.1$; $+161.8$] & $+5.8$ [$+3.1$; $+35.6$] \\
			Carta           & $+0.5$ & $+35.6$ [$+21.8$; $+50.0$] & $+34.3$ [$+2.3$; $+173.9$] & $+30.5$ [$+19.3$; $+76.6$] \\
			& $-0.5$ & $-33.7$ [$-50.0$; $-15.7$] & $+102.0$ [$+68.6$; $+739.3$] & $+29.5$ [$+6.2$; $+172.2$] \\
			Vetro           & $+0.5$ & $+100.0$ [$+50.0$; $+100.0$] & $-7.5$ [$-51.5$; $+352.9$] & $+49.6$ [$+9.2$; $+155.5$] \\
			& $-0.5$ & $-44.4$ [$-50.0$; $-38.9$] & $+682.5$ [$+194.0$; $+1171.0$] & $+38.2$ [$+5.1$; $+71.2$] \\
			\hline
	\end{tabular}}
\end{table}

Nel regime a tempo il costo è proporzionale al numero di giri, quindi mezzo ritiro in più a partire da $X^*$ aumenta il costo di $0.5/X^*$: il $20\%$ da $2.5$, il $33\%$ da $1.5$, il $100\%$ da $0.5$, come in tabella. Il disservizio, invece, si comporta in modo diverso per i diversi rifiuti. Per carta, indifferenziata e plastica mezzo ritiro in più aumenta in mediana anche il disservizio ($+12$--$34\%$): in almeno metà delle istanze la frequenza scelta a $\gamma=0.5$ è già quella di minimo disservizio (Sezione~\ref{sse:t1-quantile}), per cui il passo peggiora entrambe le componenti e $F_{\mathrm{tot}}$ cresce del $13$--$31\%$. L'organico, che a $\gamma=0.5$ sta sotto la propria frequenza di minimo disservizio, è l'unico rifiuto con un vero compromesso: mezzo ritiro in più costa il $18.7\%$ in più e riduce il disservizio del $14.4\%$, con $F_{\mathrm{tot}}$ che sale solo del $2.8\%$.

Un $F_{\mathrm{tot}}$ quasi piatto vicino all'ottimo non indica quindi una soluzione fisicamente stabile, ma una quasi indifferenza secondo i pesi: nell'esempio dell'organico, due soluzioni con valori di $F_{\mathrm{tot}}$ distanti meno del $3\%$ differiscono di circa un quinto nel costo e di un settimo nel disservizio. La sensibilità è inoltre asimmetrica: togliere mezzo ritiro aumenta il disservizio del $50$--$102\%$ (escluso il vetro), mentre aggiungerlo lo riduce al più del $14\%$ (organico), coerentemente con una penalità di sotto-servizio cinque volte quella di sovra-servizio ($\alpha/\beta=5$).

\subsection{Confronto tra algoritmi a pari frequenza}\label{sse:t2-algoritmi}

\begin{figure}[htbp]
	\centering
	\includegraphics[width=\textwidth]{immagini/fig_t2_algoritmi.png}
	\caption{CW rispetto a Greedy a pari $X_r$ (somma sui rifiuti, media sui seed): vantaggio logistico, cioè costo di Greedy diviso costo di CW meno uno, e variazione percentuale dei camion per giro. Linea: mediana sulle 23 istanze; fascia: minimo e massimo.}
	\label{fig:t2_algoritmi}
\end{figure}

La Figura~\ref{fig:t2_algoritmi} riassume il confronto tra gli algoritmi lungo tutta la griglia. Il vantaggio logistico di CW è in mediana del $10.1\%$ a $X_r=0.5$ ($4.1$--$21.5\%$ a seconda dell'istanza), dell'$11.3\%$ a $X_r=1.0$ e del $14.6\%$ da $X_r=2$ in poi ($7.2$--$39.1\%$). I camion per giro seguono il regime. A $X_r=0.5$, dove decide la capacità, CW ne usa in mediana l'$1.8\%$ in più di Greedy (da $-11.7\%$ a $+6.6\%$): il suo vantaggio sotto capacità viene solo dai chilometri. Da $X_r=2$, dove decide il turno, ne usa il $5.8\%$ in meno (fino al $-35\%$), perché il tempo risparmiato in viaggio lascia spazio a più utenti per camion. Da $X_r=2$ in poi entrambe le curve sono costanti, come atteso dal fatto che oltre la soglia il giro non cambia. Il vantaggio maggiore nel regime a tempo è il motivo per cui il gap tra gli algoritmi è più alto quando la soluzione scelta sta sopra la soglia, come a Fabriano a $\gamma=0.5$ rispetto a $\gamma=0.1$ (Sezione~\ref{se:b4}).

\subsection{Ricerca potata della griglia}\label{sse:t2-potatura}

La struttura del disservizio suggerisce una riduzione della ricerca. $F_{\mathrm{insoddis}}(X_r)$ dipende solo da $X_r$ e dalle frequenze ideali degli utenti, quindi si calcola su tutta la griglia senza costruire alcuna rotta. Sia $X^{\min}_r$ il punto della griglia con disservizio minimo, cioè la frequenza al quantile $83.3\%$ della Sezione~\ref{sse:t1-quantile}. Per $X_r>X^{\min}_r$ il disservizio non diminuisce; se nemmeno il costo logistico diminuisce, ogni $X_r>X^{\min}_r$ è peggiore di $X^{\min}_r$ per qualunque $\gamma$ e $k_{\mathrm{scala}}$ e non va instradata. La ricerca potata valuta quindi solo $X_r\le X^{\min}_r$.

La condizione sul costo non è garantita a priori, perché le rotte sono euristiche, ma è verificata in tutti i $6610$ profili dei benchmark: il costo per $X_r>X^{\min}_r$ non scende mai sotto quello di $X^{\min}_r$, e su $125\,590$ scelte ($6610$ profili per 19 valori di $\gamma$) la frequenza ottima sta sempre entro $X^{\min}_r$. La potatura valuta in media $4.78$ frequenze su 12 ($-60.1\%$): $1.9$ per il vetro, $3.0$ per la carta, $5.1$ per la plastica, $6.0$ per l'indifferenziata e $7.9$ per l'organico. Rispetto a una ricerca che parte da $X_r=0.5$ e si ferma al primo aumento di $F_{\mathrm{tot}}$, che trova il minimo nel $99.9\%$ dei casi, la potatura è esatta nei casi esaminati, non dipende da $\gamma$ e, per coprire tutti i 19 valori di $\gamma$, richiede meno valutazioni ($4.78$ contro $5.72$; per un solo valore di $\gamma$ la discesa ne richiede in media $4.60$).

La regola è stata implementata e verificata ricalcolando 29 istanze (sintetico originale con $N=100$, $500$, $1000$; scenari residenziale, villette e grandi condomini; grafo standard; sottocampioni di Fabriano con $N=1000$ e $2000$; Fabriano completa). Per ogni istanza e algoritmo sono state eseguite sia la griglia completa sia quella potata, con lo stesso insieme di processi paralleli. Le $1382$ righe della griglia potata sono identiche a quelle della griglia completa e le frequenze scelte coincidono in tutte le $5510$ scelte (29 istanze, 2 algoritmi, 5 rifiuti, 19 valori di $\gamma$). Anche le $4862$ righe ricalcolate coincidono con quelle dei CSV dei benchmark, entro la precisione con cui questi sono salvati: i risultati sono riproducibili.

\newpage
\begin{table}[htbp]
	\centering
	\caption{Ricerca potata contro griglia completa (stesso insieme di processi, stessa macchina; mediana sui seed): quota di frequenze valutate, rapporto tra i tempi reali e tempo reale della griglia completa.}
	\label{tab:t2_tempi}
	\small
	\begin{tabular}{lccccc}
		\hline
		& & \multicolumn{2}{c}{Tempo potata / completa} & \multicolumn{2}{c}{Completa (s)} \\
		Istanza & Valutazioni & CW & Greedy & CW & Greedy \\
		\hline
		Sintetico, $N=100$         & 0.40 & 0.86 & 0.85 & 5.4  & 5.4 \\
		Sintetico, $N=500$         & 0.40 & 0.80 & 0.84 & 6.1  & 5.6 \\
		Sintetico, $N=1000$        & 0.40 & 0.68 & 0.80 & 8.8  & 6.1 \\
		Residenziale, $N=400$      & 0.40 & 0.82 & 0.84 & 5.9  & 5.6 \\
		Villette, $N=400$          & 0.27 & 0.69 & 0.70 & 5.8  & 5.5 \\
		Grandi condomini, $N=400$  & 0.52 & 0.95 & 0.98 & 5.9  & 5.5 \\
		Grafo standard, $N=200$    & 0.40 & 0.85 & 0.86 & 5.5  & 5.4 \\
		Fabriano, $N=1000$         & 0.40 & 0.67 & 0.82 & 8.7  & 6.2 \\
		Fabriano, $N=2000$         & 0.40 & 0.55 & 0.72 & 21.7 & 7.8 \\
		Fabriano, $N=3932$         & 0.40 & 0.43 & 0.60 & 87.6 & 16.8 \\
		\hline
	\end{tabular}
\end{table}

Il risparmio di tempo è minore di quello di valutazioni e cresce con la dimensione dell'istanza (Tabella~\ref{tab:t2_tempi}). Fino a $N=500$ la griglia completa richiede tra $5.4$ e $6.1$\,s indipendentemente da $N$, segno di un costo fisso (avvio dei processi e trasferimento dei dati) che la potatura non riduce, e il tempo scende solo del $14$--$20\%$. Su Fabriano completa, dove prevale il routing, il tempo di CW passa da $87.6$ a $37.7$\,s ($-57\%$, vicino al $-60\%$ delle valutazioni) e quello di Greedy da $16.8$ a $10.1$\,s ($-40\%$). Sulle 29 istanze il tempo totale scende del $29.7\%$ e il tempo di CPU del $31.9\%$. Il numero di frequenze valutate dipende dal mix tipologico, perché questo fissa il quantile: $27\%$ della griglia nelle villette, $52\%$ nei grandi condomini.

L'applicazione operativa ha due limiti. Il primo è che la condizione sul costo è verificata empiricamente e non garantita: un'implementazione prudente dovrebbe valutare anche il primo punto oltre $X^{\min}_r$ e controllare che il costo non scenda. Il secondo riguarda $k_{\mathrm{scala}}$, che nei benchmark è calcolato sulla griglia completa di Greedy: con la ricerca potata va definito in altro modo, oppure la scelta va espressa direttamente come prezzo accettato per unità di disservizio (Sezione~\ref{sse:t1-soglie}), che non richiede $k_{\mathrm{scala}}$.

\subsection{Sintesi}\label{sse:t2-sintesi}

A frequenza fissata la soluzione attraversa due regimi separati da una soglia prevedibile, $\bar X_r=W_r/(C_r V^{\mathrm{t}}_r)$, che coincide con quella osservata nel $95.8\%$ dei profili e non la anticipa mai. Oltre la soglia il giro è fisso, il costo cresce in proporzione al numero di giri e il vantaggio di CW si stabilizza al $14.6\%$ in mediana, con il $5.8\%$ di camion in meno; sotto la soglia il vantaggio è minore e non viene dai camion. Vicino all'ottimo un $F_{\mathrm{tot}}$ piatto può nascondere differenze di un quinto nel costo, e togliere mezzo ritiro costa più disservizio di quanto aggiungerlo faccia risparmiare. Infine, poiché il disservizio non richiede il routing, la griglia si può potare alla frequenza di minimo disservizio: in tutti i casi esaminati la soluzione non cambia, le valutazioni scendono del $60\%$ e il tempo di calcolo su Fabriano completa del $40$--$57\%$.

\section{Discussione complessiva}\label{se:discussione}

Questa sezione riunisce i risultati dei cinque benchmark e dei due benchmark trasversali in quattro domande: che cosa decide $\gamma$, quanto vale Clarke-Wright rispetto a Greedy e perché, che cosa cambia da un'istanza all'altra e quanto sono solide queste conclusioni. I valori riportati sono quelli delle sezioni precedenti, a cui si rimanda per i dettagli.

\subsection{Che cosa decide \texorpdfstring{$\gamma$}{gamma}}\label{sse:disc-gamma}

Il parametro $\gamma$ non modifica le rotte: a frequenza fissata camion, chilometri e ore sono gli stessi per ogni $\gamma$, che interviene solo nella scelta di $X_r^*$ (Sezione~\ref{se:metodo-analisi}). Per questo la soluzione è costante a tratti: sui 19 valori di $\gamma$ le soluzioni distinte sono da $6$ a $10$ per istanza e la stessa soluzione resta ottima, in mediana, per $\gamma$ tra $0.55$ e $0.90$ (Sezione~\ref{sse:t1-soglie}). Il confronto tra politiche va quindi fatto sulla soluzione fisica e non sul valore di $F_{\mathrm{tot}}$, che cambia anche solo perché cambiano i pesi.

In unità fisiche, passare da $\gamma=0.10$ a $\gamma=0.50$ riduce il disservizio del $77\%$ (mediana sulle istanze, CW) e aumenta del $157$--$175\%$ costo, chilometri e ore di lavoro; passare da $0.50$ a $0.90$ toglie altri $4$ punti di disservizio al prezzo di altri $36$ punti di costo (Tabella~\ref{tab:t1_fisica}). Il meccanismo fisico è il cambio di regime (Sezione~\ref{sse:t2-regimi}). Sotto la soglia $\bar X_r=W_r/(C_rV^{\mathrm{t}}_r)$ il carico per ritiro rende vincolante la capacità, sopra quella soglia il turno: il giro non cambia più e il costo cresce in proporzione al numero di giri. Ne deriva il paradosso dei camion: la politica più economica usa più camion per giro (a Fabriano $145$ contro $93$ per Greedy), perché con frequenze basse ogni camion si riempie prima di esaurire il turno, mentre le attivazioni settimanali e il lavoro sono minimi. L'effetto cresce con il carico per fermata, da assente nelle villette a un rapporto di $3.3$ nei grandi condomini (Sezione~\ref{se:b2}). Tra $\gamma=0.50$ e $\gamma=0.90$, invece, la flotta per giro non cambia mai: cambia solo quante volte la si usa.

Il disservizio non si annulla mai. Con una frequenza unica per rifiuto e frequenze ideali diverse tra tipologie, il minimo raggiungibile cade al quantile $\alpha/(\alpha+\beta)=83.3\%$ delle frequenze ideali degli utenti, previsto esattamente in tutti i 30 casi verificati (Sezione~\ref{sse:t1-quantile}). Avvicinandosi a quel quantile ogni mezzo ritiro in più evita sempre meno disservizio, mentre nel regime a tempo costa sempre lo stesso: il costo marginale cresce quindi a ogni passo, in tutte le coppie istanza-seed, da $0.087$\,\euro\ a $4.95$\,\euro\ per unità di disservizio evitata (mediane del primo e dell'ultimo passo). Il salto finale nasce dal fatto che il quantile cade a metà di una cella della griglia: a Fabriano l'ultimo passo costa $7.9$\,\euro\ per unità, aumenta il costo logistico dell'$11.6\%$ e riduce il disservizio dell'$1.8\%$. Nelle villette e nei grandi condomini, dove il quantile cade su un punto della griglia, il salto finale non esiste.

Per il decisore ne seguono due indicazioni. La prima è che la scelta è tra poche soluzioni discrete, ciascuna con un prezzo esplicito per unità di disservizio evitata: la soluzione del plateau è l'ultima prima del salto, e la successiva va giustificata con quel prezzo. La seconda è che $\gamma$ non è trasferibile da un'istanza all'altra, perché un passo di costo $s$ entra a $\gamma^*=s/(k_{\mathrm{scala}}+s)$ e $k_{\mathrm{scala}}$ varia tra $0.55$ e $2.05$: la grandezza da fissare è il prezzo accettato, non il peso.

\subsection{Quanto vale Clarke-Wright e perché}\label{sse:disc-algoritmi}

Il confronto corretto tra i due algoritmi è a parità di frequenze, cioè a parità di disservizio. In questo confronto CW costa meno di Greedy nel $95.3\%$ delle combinazioni a $X_r=0.5$ e nel $99.2$--$99.3\%$ per $X_r\ge1.5$; servendo con CW le frequenze scelte da Greedy il costo scende nel $99.6$--$100\%$ dei casi, del $9$--$13\%$ in mediana (Sezione~\ref{sse:t1-algoritmi}), e a Fabriano del $6.8$--$7.0\%$ con $5$ camion in meno su $93$. La frontiera di CW sta quindi sotto quella di Greedy.

Il vantaggio ha due componenti con comportamento diverso (Sezione~\ref{sse:t2-algoritmi}). Quella dei chilometri è sempre presente: CW percorre per giro dal $30\%$ al $42\%$ di chilometri in meno nel sintetico al crescere di $N$ e il $47\%$ in meno a Fabriano. Quella dei camion esiste solo nel regime a tempo, dove il tempo risparmiato in viaggio lascia spazio a più utenti per camion ($-5.8\%$ di camion per giro in mediana); sotto capacità CW non risparmia camion e anzi ne usa in mediana l'$1.8\%$ in più, e raggiunge la soglia di regime più tardi di Greedy (nel $10.7\%$ dei profili dell'organico contro il $5.0\%$). Il vantaggio logistico passa così dal $10.1\%$ a $X_r=0.5$ al $14.6\%$ da $X_r=2$ in poi.

A pari $\gamma$, invece, i due algoritmi possono scegliere frequenze diverse. Poiché ogni aumento di frequenza costa meno a CW, la sua soglia $\gamma^*$ è più bassa nel $98.6\%$ dei passi e CW alza le frequenze prima di Greedy: per $\gamma$ tra $0.10$ e $0.20$ e tra $0.45$ e $0.55$ le soluzioni scelte non sono confrontabili fino al $58.5\%$ dei casi, perché CW ha meno disservizio ma costa di più. È il caso di Fabriano a $\gamma=0.10$, dove CW costa il $7.2\%$ in più e ha il $17.8\%$ di disservizio in meno: non è un peggioramento del routing ma una scelta di servizio diversa.

Il gap su $F_{\mathrm{tot}}$ combina questi elementi. È il prodotto del vantaggio logistico $L_g$ per il peso $w$ del costo nel totale, più lo spostamento di $X^*$, e il suo massimo cade a $\gamma=0.5$ in tutti i benchmark (nel grafo standard non distinguibile da $\gamma=0.1$) perché lì $X^*$ è già nel regime a tempo, dove $L_g$ è più alto, e il peso del costo è ancora alto. A $\gamma=0.1$ la soluzione è sotto la soglia e $L_g$ è minore, a $\gamma=0.9$ il termine di disservizio, uguale per i due algoritmi, diluisce il vantaggio: a Fabriano $L_g$ vale $5.3$, $7.5$ e $7.4\%$ e $w$ vale $0.55$, $0.556$ e $0.197$, per un gap di $3.88$, $4.21$ e $1.46\%$ (Sezione~\ref{se:b4}).

Il prezzo di CW è il tempo di calcolo, che oltre la componente fissa cresce circa come $N^2$ e più rapidamente di quello di Greedy (Sezione~\ref{se:b1}). Su Fabriano completa una griglia richiede $87.6$\,s contro $16.8$\,s misurando entrambi con un unico insieme di processi (Sezione~\ref{sse:t2-potatura}); nella campagna principale, dove ogni rifiuto avvia i propri processi, i tempi erano $100.8$ e $42.3$\,s. La ricerca potata porta CW a $37.7$\,s senza cambiare la soluzione.

\newpage
\subsection{Che cosa cambia da un'istanza all'altra}\label{sse:disc-istanze}

Le caratteristiche dell'istanza spostano molto i costi e poco il vantaggio relativo di CW, che dipende soprattutto dal peso del viaggio nel costo.

\begin{itemize}
	\item \textbf{Dimensione.} Il gap a $\gamma=0.5$ scende da $21.7\%$ a $N=100$ a una fascia tra $8.7\%$ e $10.7\%$ per $N=500$--$1000$, con un minimo a $500$ e una risalita a $750$ legati al numero intero di camion risparmiati (Sezione~\ref{se:b1}). Il risparmio relativo di chilometri di CW cresce con $N$, ma il peso del viaggio cala (dal $59.8\%$ al $28.0\%$ del tempo di giro) e i due effetti quasi si compensano.
	\item \textbf{Tipologia urbanistica.} Il mix cambia il carico per fermata, il regime e la frequenza di minimo disservizio, e quindi i costi: a $\gamma=0.5$ il costo logistico di CW è $3.1$ volte più alto nei grandi condomini che nelle villette. Il gap, invece, resta tra $11.6\%$ e $12.4\%$ in quattro scenari su cinque ($9.7\%$ nel misto periferia).
	\item \textbf{Distribuzione spaziale.} Concentrare gli utenti attorno al deposito riduce il costo del $20\%$ senza cambiare le frequenze, e riduce il gap a $\gamma=0.5$ da $11.6\%$ a $8.9\%$ perché il viaggio pesa meno ($27.8\%$ contro $39.6\%$ del tempo di giro di CW).
	\item \textbf{Rete reale.} Il gap di Fabriano ($4.2\%$ a $\gamma=0.5$) è molto più basso di quello del sintetico originale, ma non per rotte peggiori: a Fabriano il viaggio è il $6.4\%$ del tempo di giro di CW e CW percorre il $47\%$ di chilometri in meno, come nei sintetici alla stessa dimensione ($46$--$49\%$). Un sintetico con area e mix di Fabriano dà a $N=3932$ un gap del $4.58\%$, contro $8.75\%$ con l'area originale (Sezione~\ref{sse:b4-sintetico-reale}): la differenza dipende da area, densità e mix, e la struttura della rete stradale incide al più per $2.1$ punti.
\end{itemize}

\subsection{Solidità delle conclusioni}\label{sse:disc-solidita}

Le conclusioni sui benchmark sintetici si basano su 30 seed per istanza, e la variabilità tra istanze non è trascurabile: con il solo seed 42 il gap del grafo standard a $\gamma=0.1$ sarebbe stato $3.12\%$ contro una media di $8.22\%$ (Sezione~\ref{se:b5}). Fabriano completa è una singola istanza deterministica; la sua variabilità è stimata con sottocampioni di utenti reali, i cui intervalli di confidenza sottostimano l'incertezza per $N\ge2500$ perché i campioni condividono gran parte degli utenti. Le proprietà strutturali (soluzione costante a tratti, soglia di regime, frequenza di minimo disservizio, dominanza di CW a pari frequenze) valgono su tutte le 23 istanze esaminate e, dove derivano dalla formulazione, sono state ricavate analiticamente e verificate sui dati. Infine, $29$ istanze ricalcolate riproducono i CSV dei benchmark riga per riga (Sezione~\ref{sse:t2-potatura}). Restano non spiegati, e sono dichiarati come tali, il peggior impacchettamento di CW sotto capacità, la risalita del gap a $N=750$ e lo scarto positivo dei chilometri reali rispetto al sintetico a $N\le250$.